% concatenated from Discrete_Hyperbolic_Secant_Distributions_10_09_2026.tex
\documentclass[12pt]{amsart}

\usepackage{amsmath, amssymb}
\usepackage{mathtools}
\usepackage{bbm}
\usepackage[utf8]{inputenc}
\usepackage{calrsfs}
\DeclareMathAlphabet{\pazocal}{OMS}{zplm}{m}{n}
\usepackage{color}
\usepackage{overpic}
\usepackage{hyperref}
\usepackage{soul}
\usepackage{dsfont}
\usepackage{enumitem}

\usepackage[english]{babel}

\usepackage{array}

\usepackage{caption}
\captionsetup[table]{width=\textwidth}

\makeatletter
\@namedef{subjclassname@2020}{\textup{2020} Mathematics Subject Classification}
\makeatother

\usepackage{float}

\setlength{\parindent}{0em} 
\hoffset=-0.5in
\textwidth=6in

\newtheorem{theorem}{Theorem}[section]
\newtheorem{lemma}[theorem]{Lemma}
\newtheorem{prop}[theorem]{Proposition}
\newtheorem{cor}[theorem]{Corollary}
\newtheorem{conj}[theorem]{Conjecture}
\newtheorem{example}[theorem]{Example}

\theoremstyle{definition}
\newtheorem{defi}[theorem]{Definition}
\newtheorem{bsp}[theorem]{Example}

\theoremstyle{remark}
\newtheorem{remark}[theorem]{Remark}

\numberwithin{equation}{section}

\newcommand{\R}{\mathbb R}
\newcommand{\C}{\mathbb C}
\newcommand{\Q}{\mathbb Q}
\newcommand{\Rd}{{\mathbb R^d}}
\newcommand{\T}{\mathbb{T}}
\newcommand{\1}{\mathbbm{1}}
\newcommand{\E}{{\mathbb E}}
\newcommand{\Pb}{\mathbb P}
\newcommand{\N}{{\mathbb N}}
\newcommand{\A}{\mathbb A}
\newcommand{\Z}{{\mathbb Z}}
\newcommand{\m}{\text{-}}
\newcommand{\eps}{\varepsilon}
\newcommand{\ph}{\varphi}
\newcommand{\rh}{\varrho}
\newcommand{\thet}{\vartheta}
\newcommand{\eqfd}{\overset{\mathrm{fd}}{=}}
\newcommand{\eqd}{\overset{\mathrm{d}}{=}}

\DeclareMathOperator{\PV}{P.\!\!V.\!\!}
\DeclareMathOperator{\Ca}{Cap}
\DeclareMathOperator{\Si}{Si}
\DeclareMathOperator{\sech}{sech}
\DeclareMathOperator{\Res}{Res}

\allowdisplaybreaks

\begin{document}

\title[Discrete Hyperbolic Secant Distributions]{Discrete Hyperbolic Secant Distributions}

%\author{Peter Kern}
%\address{Peter Kern, Mathematical Institute, Faculty of Mathematics and Natural Sciences, Heinrich Heine University D\"usseldorf, Universit\"atsstr. 1, D-40225 D\"usseldorf, Germany}
%\email{kern@hhu.de}

%\keywords{Hyperbolic secant distribution, discrete distribution, Ramanujan series, Gaussian constant, lemniscate constant}
%\subjclass[2020]{Primary 60G51; Secondary 28A78, 28A80, 60G17, 60G18, 60G52.}

\author{Leonard J. E. Pleschberger}
\address{Leonard Jobst Eberhard Pleschberger, Neptunstra\ss e 18, D\"usseldorf D-40223 \& Sechshauser Stra\ss e 24/1/15, A-1150 Vienna.}
\email{leonard.pleschberger@gmail.com}
\keywords{Discrete hyperbolic secant distribution, elliptic lambda function, Ramanujan's notebooks, lemniscate constant, Gau\ss' constant, Poisson summation, Mellin transform, Riemann zeta function}


\date{\today}

\maketitle


\begin{abstract}
We introduce a family of discrete hyperbolic secant distributions on $\mathbb{Z}$, whose normalizing constants arise from series values calculated by Ramanujan and are expressed in terms of Gau\ss' constant $G=\varpi/\pi$, where $\varpi=\Gamma^2(1/4)/(2\sqrt{2\pi})$ is the lemniscate constant. Using the elliptic lambda-star function $\lambda^*$, we construct scaled versions of these distributions parametrized by $\sqrt{r}$ and $1/\sqrt{r}$ for $r \in \mathbb{N}$. For the first such distribution, we compute the moments up to the eighth degree in closed form via Poisson summation, exploiting that the hyperbolic secant is a fixed point of the $-2\pi i$-Fourier transform. As a byproduct, we obtain closed-form values for series of odd powers of the hyperbolic secant up to the ninth degree, e.\,g.\ $\sum_{k \in \mathbb{Z}} \operatorname{sech}^3(\pi k) = (G^3+G)/\sqrt{2}$, and reinterpret several classical Ramanujan series probabilistically. Finally, we apply our results to evaluate a contour integral, on the critical line, of the product of Dirichlet's beta, gamma, and Riemann zeta functions.
\end{abstract}

\vspace{-0.5cm}

\section{Introduction}

In Section 2, we introduce a family of discrete hyperbolic secant distributions. The normalizing constants are derived from values of series calculated by Ramanujan. Using Gau\ss' constant $G=\varpi/\pi$ with the lemniscate constant $\varpi=\Gamma^2(1/4)/(2\sqrt{2 \pi})$, we can represent these distributions in very basic terms, e.\,g.
$$
\left(\mathbb{P} (X = k)\right)_{k \in \mathbb{Z}} = \left(\frac{\operatorname{sech}\left(\pi k\right)}{\sqrt{2}G} \right)_{k\in\mathbb{Z}}
$$
and versions with a scaling factor of $\sqrt{r}$ and $1/\sqrt{r}$ applied to $\pi k$. We use Ramanujan series and apply the lambda-star function $\lambda^*$ which yields the required singular values of the occurring elliptic integrals.\\

In Section 3, we calculate the moments up to the eighth degree for the first hyperbolic secant distribution, e.\,g. $\mathbb{E}\left[X^2 \right] = \operatorname{Var}(X) = (G/2)^2$. Therefrom, we derive the values of series of odd multiplicities of the hyperbolic secant up to the ninth degree, e.\,g.
$$
\sum_{k\in\mathbb{Z}} \mathrm{sech}^3 (\pi k) = \left(G^3 + G\right)/\sqrt{2}.
$$
In this case, the hyperbolic secant is a fix-point of the $-2\pi i$-Fourier transform which allows us comfortable calculations via Poisson summation.\\

In Section 4, we give further statistical interpretations of some Ramanujan series in terms of the second hyperbolic secant distribution, e.\,g. $\mathbb{E}\left[(2Y+1)^2 \right] = 2G^2$.\\

In Section 5, we apply the first hyperbolic secant distribution to contour integrals. We calculate a contour integral for a product $\Xi$ of Dirichlet's beta function, the gamma function and Riemann's zeta function on the critical line, viz.
$$
\int_{1/2 - i\cdot \infty}^{1/2 + i\cdot \infty} \pi^{-(s+2)} \beta(s+2) \Gamma(s+2) \zeta(s) \, \mathrm{d}s = i \left( \frac{\varpi}{\sqrt{2}} \mathbb{E}\left[X^2 \right] - \frac{\pi}{8} \right) = i \left( \frac{\varpi}{\sqrt{2}} \left(\frac{G}{2}\right)^2 - \frac{\pi}{8} \right).
$$
We use the Mellin transform and the residue theorem.


% \section{On Bernoulli's Lemniscate}

% We explore one of the fundamental objects in planar Euklidean geometry: The \textit{lemniscate of Bernoulli}.\footnote{Jakob Bernoulli introduced his lemniscate in 1694. The word ``lemniskate'' comes from the ancient Greek substantive
% \begin{ibycus}
% lhmni'skos
% \end{ibycus}(``woolen ribbon''); hence a spelling with ``k'' instead of ``c'' should be considered. The word was retained in the Latin adjective \textit{lemniscatus} ("decorated with dangling ribbons") which derivates from the substantive \textit{lemniscus} (originally a ``ribbon of linden bast'', later crafted from precious metals as a symbol of honour). With Latin as the lingua franca of scholarship, the customary spelling with ``c'' became established.}
% The lemniskate is a closed curve shaped like the symbol of infinity. It consists of two symmetric loops with exactly one double point. Any point of the lemniskate of Bernoulli can be constructed by straightedge and compass given just two distinct points $F_1, F_2 \in \mathbb{R}^2$, call them \textit{foci}, as follows: Draw a square $\square \subset \mathbb{R}^2$ with two opposite vertices $F_1, F_2$ and centre $M$. Consider the Euklidean \textit{focal distance} $c = d(F_1, M)$. Pythagoras' theorem yields the side-length $a = \sqrt{2}c$ of $\square$. Draw two circles $S_1, S_2$ with radius $a$ around $F_1, F_2$. The lemniscate of Bernoulli is the locus of midpoints of line segments whose ends lie in $S_1$ and $S_2$ at distance $2c$. Usually, its scaling is parametrized either by its focal distance $c=d(F_1, F_2)/2 >0$ or by its half-perimeter $a = \sqrt{2}c >0$.

% For ease, we only consider those lemniscates with foci $F_1, F_2$ on the $x$-axis and double point at the origin $O$. We can describe them algebraically as the quartic
% $$
% L_c := \left\{(x,y)^\top \in \mathbb{R}^2 : (x^2 + y^2)^2 = a^2 (x^2 - y^2) = 2 c^2 (x^2-y^2)\right\}.
% $$
% One may ask: What is the most natural choice for $c$? By choosing $c=1/\sqrt{2}$, the scaling parameter on the right hand side of the quartic above becomes $1$, 



% Gau\ss \ chose $c_1$ in such a way that the right loop of the lemniscate $L_{c_1}$, i.~e. the intersection of $L_{c_1}$ with the closed first and fourth quadrant, has an arc-length which equals his \textit{lemniscate constant}
% $$
% \varpi := \frac{\Gamma^2(1/4)}{2\sqrt{2\pi}}.
% $$
% This constant determines the scaling of Gau\ss' \textit{unit lemniscate}, i.e. the focal distance $c_1$:\\

% \textbf{Theorem.} Let the right loop of Gau\ss' unit lemniscate $L_{c_1}$ have an arc-length of $\varpi$. This implies a focal distance of
% $$
% c_1 = \frac{1}{\sqrt{2}}.
% $$

% \textbf{Proof:} We calculate the arc-length of the right loop of $L_{c_1}$ by switching to polar coordinates. For $r > 0$ and $\varphi \in [0,2\pi)$, let $x = r(\varphi) \cos(\varphi)$ and $y = r(\varphi) \sin(\varphi)$. Further, Pythagoras's theorem applied to the unit circle yields the trigonometric identity $\sin^2(\varphi) + \cos^2(\varphi) = 1$ and we use the double-angle formula $\cos(2\varphi) = \cos^2(\varphi) - \sin^2(\varphi)$. Our parametrization of the lemniscate $L_{c_1}$ yields
% $$
% (x^2 + y^2)^2 = 2 c_1^2 \cdot (x^2-y^2) \ \Leftrightarrow \ r^2 = 2c_1^2 \cdot \cos(2\varphi).
% $$
% Since $r^2, c_1 > 0$, we must have $\cos(2\varphi) > 0$ which is the case for $\varphi \in [0,\pi/4)$ in the first quadrant. Due to its symmetry, the perimeter $U_{L_c}^+$ of the right loop is obtained by doubling the arc-length of its restriction to the first quadrant. Hence, in terms of $x(\varphi), \ y(\varphi)$, this arc-length is obtained via the substitution $2\varphi = \vartheta \ \Rightarrow \ \mathrm{d}\varphi = \mathrm{d}\vartheta/2$ as follows:
% \begin{align*}
%     U_{L_c}^+ &= 2 \int_0^{\pi/4} \sqrt{(x'(\varphi))^2 + (y'(\varphi))^2} \, \mathrm{d}\varphi\\
%     &= 2 \int_0^{\pi/4} \sqrt{(r'(\varphi)\cos(\varphi) + r(\varphi)\sin(\varphi))^2 + (r'(\varphi)\sin(\varphi) + r(\varphi)\cos(\varphi))^2} \, \mathrm{d}\varphi\\
%     &= 2 \int_0^{\pi/4} \sqrt{r^2(\varphi) + \left(\frac{\mathrm{d}}{\mathrm{d}\varphi}\left[\sqrt{2c_1^2\cos(2\varphi)} \right] \right)^2} \, \mathrm{d}\varphi\\
%     &= 2\sqrt{2}c_1 \int_0^{\pi/4} \frac{1}{\sqrt{\cos(2\varphi)}} \, \mathrm{d}\varphi\\
%     &= \sqrt{2}c_1 \int_0^{\pi/2} \frac{1}{\sqrt{\cos(\vartheta)}} \, \mathrm{d}\vartheta.
% \end{align*}
% This is not solvable by elementary analysis, hence we have to enter the realm of special functions, cf. "Handbook of Mathematical Functions with Formulas, Graphs, and Mathematical Tables" from Abramowitz and Stegun. We substitute 
% $$
% t = \sin^2(\vartheta) \ \Leftrightarrow \ \cos(\vartheta) = \sqrt{1 - t}
% $$ 
% which implies 
% $$\mathrm{d}t = 2\sin(\vartheta)\cos(\vartheta)\, \mathrm{d}\vartheta = 2 \sqrt{t} \cdot \sqrt{1-t} \, \mathrm{d}\vartheta \ \Leftrightarrow \ \mathrm{d}\vartheta = \mathrm{d}t/(2 \sqrt{t} \cdot \sqrt{1-t}).
% $$
% By using the \textit{Beta function} $B$, also called \textit{Euler's integral of the first kind}, viz.
% $$
% B(z_1, z_2) := \int_0^1 t^{z_1-1} (1-t)^{z_2-1} \, \mathrm{d}t
% $$
% for complex numbers $z_1,z_2 \in \mathbb{C}$ with $\Re(z_1), \Re(z_2) > 0$
% and further the \textit{Gamma function} 
% $$
% \Gamma(z) := \int_0^1 t^{z-1} e^{-t} \, \mathrm{d}t
% $$
% for complex numbers $z \in \mathbb{C}$ with $\Re(z) > 0$, we get
% $$
% U_{L_c}^+ = \frac{\sqrt{2}c_1}{2} \int_0^{1} t^{-1/2} (1 - t)^{-3/4} \, \mathrm{d}t = \frac{\sqrt{2}c_1}{2} \cdot B(1/2, 1/4) = \frac{\sqrt{2}c_1}{2} \cdot  \frac{\Gamma(1/2) \cdot \Gamma(1/4)}{\Gamma(3/4)}.
% $$
% Now, Euler's reflection formula $\Gamma(z) \cdot \Gamma(1-z) = \pi \cdot \csc(\pi z)$ yields $\Gamma(3/4) = \sqrt{2}\pi / \Gamma(1/4)$. With the identity $\Gamma(1/2) = \sqrt{\pi}$ we get
% $$
% U_{L_c}^+ = \sqrt{2}c_1 \cdot \frac{\Gamma^2(1/4)}{2\sqrt{2\pi}} = \sqrt{2}c_1 \cdot \varpi = \varpi \ \Leftrightarrow \ c_1 = \frac{1}{\sqrt{2}}
% $$
% for the lemniscate constant $\varpi$. \hfill $\square$\\

% \textbf{Total Arc-length:} The proof of the previous Theorem directly yields the following generalization:\\

% \textbf{Corollary:} For $a/\sqrt{2} = c > 0$, let $L_c$ be a lemniscate of Bernoulli. Then its right loop has the arc-length
% $$
% U_{L_c}^+ = \sqrt{2}c \cdot \varpi = a \cdot \varpi,
% $$
% with the lemniscate constant $\varpi = \Gamma^2(1/4)/(2\sqrt{2\pi})$. This yields the arc-length of the whole lemniscate, i.~e.
% $$
% U_{L_c} = 2\sqrt{2}c \cdot \varpi = 2a \cdot \varpi
% $$
% due to symmetry.
% \\

% \textbf{Proportionality Constant:} Gau\ss \ defined the lemniscate constant as the proportion of the arc-length of the right loop of a lemniscate to the radius $r=a$ of a circumcircle $\bigcirc_a$ of $L_{c}$ with $a=\sqrt{c} > 0$, i.e.
% $$
% \frac{U_{L_c}^+}{a} = \varpi = \frac{\Gamma^2(1/4)}{2\sqrt{2\pi}}.
% $$
% In analogy to $\pi$, one could also take the proportion between the arc-length of the whole lemniscate and the diameter $d=2a$ of the circumcircle $\bigcirc_a$ of a lemniscate $L_c$ with $a = \sqrt{c} > 0$. This yields
% $$
% \frac{U_{L_c}}{d} = \varpi = \frac{\Gamma^2(1/4)}{2\sqrt{2\pi}}.
% $$
% A third alternative yields a probably more elegant doubled lemniscate constant $\widetilde{\varpi}$: One could take the proportion between the arc-length of the whole lemniscate and the radius $r=a$ of the circumcircle $\bigcirc_a$ of a lemniscate $L_c$ with $a = \sqrt{c} > 0$. This would yield
% $$
% \frac{U_{L_c}}{a} = 2 \cdot \varpi = \widetilde{\varpi} = \frac{\Gamma^2(1/4)}{\sqrt{2\pi}}.
% $$


% \textbf{Total Area:} We also derive a formula for the area enclosed by the right loop of $L_c$:\\

% \textbf{Proposition:} For $a / \sqrt{2} = c > 0$, let $L_c$ be a lemniscate of Bernoulli. Then the area enclosed by its right loop, name it $A_{L_c}^+$ has the Lebesgue measure
% $$
% \lambda\left(A_{L_c}^+\right) = \frac{a^2}{2} = c^2.
% $$
% Hence, the total area enclosed by $L_c$, call it $A_{L_c}$, has Lebesgue measure
% $$
% \lambda\left(A_{L_c}\right) = a^2 = 2c^2
% $$
% due to the symmetry.\\

% \textbf{Proof:} Because of the symmetry, we calculate $2$ times the area of $A_{L_c}$ intersecting with the fist quadrant. We use the the same substitutions and formulas as for the perimeter $U$ from above and get
% $$
% \lambda\left(A_{L_{c}}\right) = 2\int_{0}^{\pi/4} \frac{r^2(\varphi)}{2} \, \mathrm{d}\varphi = 2c^2 \int_0^{\pi/4} \cos(2\varphi) \, \mathrm{d}\varphi = 2 c^2 \left[ \frac{\sin(2\varphi)}{2} \right]_{\varphi = 0}^{\pi/4} = c^2,
% $$
% as stated. \hfill $\square$\\


\section{Basic Discrete Hyperbolic Secant Distributions}

We introduce a family of six discrete hyperbolic secant distributions through their probability mass functions. These functions are normalized by series values given in \cite{ramanujan_1991} by using the lemniscate modulus $k_1 = 1/\sqrt{2}$. We represent the normalizing constants in terms of the Gau\ss ian constant $G:=\varpi/\pi$ with the lemniscate constant $\varpi:=\Gamma^2(1/4)/(2\sqrt{2 \pi})$. This results in particularly simple expressions.

\begin{theorem} \label{thm:pmfs}
We tabulate six discrete hyperbolic secant probability distributions:

\renewcommand{\arraystretch}{2}
\begin{table}[h]
    \centering
    \begin{tabular}{c|c|c|c}
PMF  &   Variable & Constant  &  Law \\
\hline
$(p_n)_{n \in \mathbb{N}_0}$ & $X_u$ & $C_{1,u}$ & $\frac{2}{\sqrt{2}G  + 1} \operatorname{sech} \left(\pi n\right)$\\
$(q_k)_{k \in \mathbb{Z}}$ & $X_b$ & $C_{1,b}$ &  $\frac{1}{\sqrt{2}G} \operatorname{sech}\left(\pi k\right)$\\
$(r_n)_{n \in \mathbb{N}_0}$ &  $Y_u$ & $C_{2,u}$ &  $\frac{2}{G} \operatorname{sech} \left(\frac{\pi}{2} (2n+1)\right)$\\
$(s_k)_{k \in \mathbb{Z}}$ &  $Y_b$ & $C_{2,b}$ &  $\frac{1}{G} \operatorname{sech}\left(\frac{\pi}{2} (2k+1)\right)$\\
$(t_n)_{n \in \mathbb{N}_0}$ & $Z_u$ & $C_{3,u}$ &  $\frac{2}{(1+\sqrt{2})G + 1} \operatorname{sech}\left(\frac{\pi}{2} n\right)$\\
$(u_k)_{k \in \mathbb{Z}}$ &  $Z_b$ & $C_{3,b}$ & $\frac{1}{(1+\sqrt{2})G}  \operatorname{sech}\left(\frac{\pi}{2} k\right)$
\end{tabular}
    % \caption{Note, that the probability mass functions for the square numbers in \textbf{bold} have quiet simple normalization constants in terms of Gau\ss' constant $G=\varpi/\pi$.}
    % \label{tab:sqrt_r_sech}
\end{table}

In full length, we name them the first, second and third discrete uni- or bilateral hyperbolic secant distributions and omit adjectives if the context is clear.
\end{theorem}

\begin{proof}
The normalizing constants can be deduced by the astonishing work of Ramanujan: He wrote in his \textit{Entry 17 (i)} in Chapter $17$ of his second notebook [in modern notation] his identity
\begin{equation}\label{form:ramanujan_1}
1 + 2 \sum_{n=1}^\infty \operatorname{sech}\left(\frac{K'(k)}{K(k)} \pi n\right) = \frac{2}{\pi} K(k), \quad \mathrm{with} \quad K(k) = \int_0^{\pi/2} \frac{\mathrm{d}\vartheta}{\sqrt{1-k^2\sin^2(\vartheta)}},
\end{equation}
Jacobi's complete elliptic integral of the first kind. We use the usual notation $K'(k) = K(k') = K(\sqrt{1-k^2})$ with the elliptic modulus $k$. The identity is proven by Berndt, cf. p. 138 of \cite{ramanujan_1991}. The $\lambda^\star$-function gives the value for the elliptic modulus $\lambda^\star(r)=k_r$ for which
\begin{equation}\label{form:elliptic_modulus}
\frac{K'(\lambda^\star(r))}{K(\lambda^\star(r))} = \sqrt{r} \quad \mathrm{with} \quad k_1 = \lambda^\star(1) = \frac{1}{\sqrt{2}}, \quad K(k_1) = \frac{\varpi}{\sqrt{2}}
\end{equation}
holds. For all the analysis, cf. the classical text  \cite{whittaker_1920}. Hence, by taking $r=1$, we get
$$
\sum_{n=0}^\infty \operatorname{sech}(\pi n) = \frac{\varpi}{\sqrt{2} \cdot \pi} + \frac{1}{2} = \frac{\sqrt{2} G + 1}{2} \ \Rightarrow \ C_{1,u} = \frac{2}{\sqrt{2}G + 1} \ \Rightarrow \ C_{1,b} = \frac{1}{\sqrt{2}G}
$$
due to the symmetry. In \textit{Entry 16 (ix)} of Chapter 17 of his second notebook, Ramanujan stated [in modern notation] the identity
$$
\sum_{n=0}^\infty \operatorname{sech}\left(\frac{\pi}{2} (2n+1)\right) = \frac{\varpi}{2 \pi} = \frac{G}{2}, \quad \Rightarrow \quad C_{2,u} = \frac{2}{G}.
$$
This is proven by Berndt on p. 134 of \cite{ramanujan_1991} and by Campbell in Theorem 7 of \cite{campbell_2023}. Fom  the absolute convergence follows
\begin{align*}
 & \sum_{n=0}^\infty \operatorname{sech}\left(\frac{\pi}{2}\cdot 2n \right) + \sum_{n=0}^\infty \operatorname{sech}\left(\frac{\pi}{2} \cdot (2n + 1)\right) = \frac{(1+\sqrt{2})G + 1}{2}.
\end{align*}
Hence, we get
$$
C_{3,u} = \frac{2}{(1+\sqrt{2})G + 1} \quad \mathrm{and} \quad C_{3,b} = \frac{1}{(1+\sqrt{2})G},
$$
where the latter was proven on p. 137 in \cite{ramanujan_1991}. Finally, 
$$
\sum_{k \in \mathbb{Z}} \operatorname{sech}\left(\frac{\pi}{2}(2k+1)\right) = \sum_{k \in \mathbb{Z}} \operatorname{sech}\left(\frac{\pi}{2} \cdot k\right) - \sum_{k\in \mathbb{Z}} \operatorname{sech} \left(\frac{\pi}{2} \cdot 2k\right) = G, \ \Rightarrow \ C_{2,b} = \frac{1}{G}
$$
finishes the proof.
\end{proof}

\vspace{1cm}

If one wanted to exaggerate, one could interpret $1 + \sqrt{2}$  as the silver ratio, but we tend to regard the expression merely as the sum of $1$ and $\sqrt{2}$.\\

The proof of the last theorem just handled the lemniscate case for $r=1$ in (\ref{form:elliptic_modulus}). In the following, we use the notations introduced in this proof. Now, let $r \ge 1$. By Ramanujan's formula \ref{form:ramanujan_1}, we immediately get a scaled family of discrete hyperbolic secant distributions.

\begin{cor}
Let $r \in \mathbb{N}$. Then, the formula
$$
\left(p_{r,k}\right)_{k\in\mathbb{Z}} = \mathbb{P}(X=k) \left(\frac{\pi}{2K(k_r)}\operatorname{sech} (\sqrt{r}\pi k)\right)_{k\in\mathbb{Z}}
$$
is the probability mass function of a random variable $X \sim \mathsf{Sech}(\sqrt{r})$. We call this distribution the discrete bilateral hyperbolic secant distribution with scaling factor $\sqrt{r}$.
\end{cor}


We tabular some interesting singular values $k_r$ with corresponding probability mass functions $(p_{r,k})_{k \in \mathbb{Z}}$

\renewcommand{\arraystretch}{2}
\begin{table}[h]
    \centering
    \begin{tabular}{c|c|r}
$p_{r,k}$  &   $ \lambda^*(r) = k_r$ & $\frac{\pi}{2K(k_r)}\operatorname{sech} (\sqrt{r}\pi k)$\\
\hline
$\mathbf{p_{1,k}}$  & $\frac{1}{\sqrt{2}}$ & $\frac{1}{\sqrt{2}G} \operatorname{sech} (\pi k)\quad \ $\\
$p_{2,k}$  & $\sqrt{2}-1$ & $\frac{4 \pi\sqrt{2\pi}}{\sqrt{\sqrt{2}+1} \cdot \Gamma(1/8)\Gamma(3/8)} \operatorname{sech} (\sqrt{2}\pi k)$\\
$p_{3,k}$  & $\frac{\sqrt{2}\left(\sqrt{3}-1\right)}{4}$ & $\frac{\sqrt[4]{2^3} \cdot \pi^2 }{\sqrt[4]{3} \cdot \Gamma^3(1/3)}\operatorname{sech} (\sqrt{3}\pi k)$\\
$\mathbf{p_{4,k}}$  & $3-2\sqrt{2}$ & $\frac{2}{(1 + \sqrt{2})G}\operatorname{sech} (2\pi k)\quad$\\
$p_{5,k}$  & $\frac{1}{\sqrt{2}} - \sqrt{\sqrt{5}-2}$ & $\frac{2\pi \sqrt{10 \pi}}{\sqrt[4]{2+\sqrt{5}} \sqrt{\Gamma(1/20)\Gamma(3/20)\Gamma(7/20)\Gamma(9/20)}} \operatorname{sech} (\sqrt{5}\pi k)$\\
$p_{6,k}$  & $\left(2-\sqrt{3}\right)\left(\sqrt{3}-\sqrt{2}\right)$ & $\frac{4\pi \sqrt{6 \pi}}{\sqrt{(1 +\sqrt{2}+\sqrt{6}) \cdot \Gamma(1/24)\Gamma(5/24)\Gamma(7/24)\Gamma(11/24)}} \operatorname{sech} (\sqrt{6}\pi k)$\\
$p_{7,k}$  & $\frac{3-\sqrt{7}}{8}$ & $\frac{2\sqrt[4]{7} \cdot \pi^2}{\Gamma(1/7)\Gamma(2/7)\Gamma(4/7)} \operatorname{sech} (\sqrt{7}\pi k)$\\
% $p_{8,k}$  & $\frac{\pi}{2K\left(\left(1 + \sqrt{2} -\sqrt{2\sqrt{2}+2}\right)^2\right)} \operatorname{sech} (\sqrt{8} \pi k)$ & $\frac{2^{1/4}\cdot 8 \pi^{3/2}}{\left(\sqrt{2}+1\right)^{1/4} \sqrt{2\sqrt{2} + \sqrt{1 + 5 \sqrt{2}}} \cdot \Gamma(1/8)\Gamma(3/8)} \operatorname{sech} (\sqrt{8}\pi k)$\\
$\mathbf{p_{9,k}}$  & $\frac{\left(\sqrt{2}-\sqrt[4]{3}\right)\left(\sqrt{3}-1\right)}{2}$ & $\frac{3}{\left(\sqrt[4]{3^3} + \sqrt[4]{3} \right) G} \operatorname{sech} (3\pi k)\quad$\\
$\mathbf{p_{16,k}}$  & $\left(1 + \sqrt{2}\right)^2 \left(\sqrt[4]{2} - 1\right)^4$ & $\frac{4}{\left(1 + \sqrt[4]{2} \right)^2 G} \operatorname{sech} (4\pi k)\quad$\\
$\mathbf{p_{25,k}}$  & $\frac{\left(3 - 2 \sqrt[4]{5}\right)\left(\sqrt{5} - 2\right)}{\sqrt{2}}$ & $\frac{5}{\left(2\sqrt{2} + \sqrt{10} \right) G} \operatorname{sech} (5\pi k) \quad$
\end{tabular}
    \caption{Note, that the probability mass functions for the square numbers in \textbf{bold} have quiet simple normalization constants in terms of Gau\ss' constant $G=\varpi/\pi$.}
    \label{tab:sqrt_r_sech}
\end{table}

We can also determine the probability mass functions for the DSH distribution with scaling factor $1/\sqrt{r}$. 
The DSH distribution makes an easy proof of the following functional equation possible, which was shown in \cite{kirschenhofer_1996} by using the Mellin transform. For $x > 0$, the functional equation
\begin{equation}\label{eq:functional_F}
F(x) := \sum_{k=1}^\infty \frac{e^{-kx}}{1 + e^{-2kx}} =  \frac{\pi}{4x} - \frac{1}{4} + \frac{\pi}{x}F\left(\frac{\pi^2}{x}\right)
\end{equation}
holds. In terms of the hyperbolic secant, we can also write 
$$
F(x) = \frac{1}{4} \sum_{k\in\mathbb{Z}} \operatorname{sech}(xk) - \frac{1}{4}
$$
and get the following scaling theorem:

\begin{theorem}
Let $x \in \mathbb{R}$. Then one has the functional equation
\begin{equation}
\sum_{k \in \mathbb{Z}} \operatorname{sech}\left(xk\right) = \frac{\pi}{x} \sum_{k \in \mathbb{Z}} \operatorname{sech}\left(\frac{\pi^2}{x} k \right)
\end{equation}
which is equivalent to (\ref{eq:functional_F}). For every $r \in \mathbb{N}$, this yields the probability mass function
$$
\left(q_{r,k} \right)_{k \in \mathbb{Z}} = \left(\frac{\pi}{\sqrt{r} \cdot 2K(k_r)} \operatorname{sech}\left( \frac{\pi}{\sqrt{r}} k \right) \right)_{k \in \mathbb{Z}}.
$$
for the discrete bilateral hyperbolic secant distribution with scaling factor $1/\sqrt{r}$.
\end{theorem}

\begin{proof}
The $-2i \pi$-Fourier transform $\mathcal{F}$ of the hyperbolic secant is well-known, viz.
$$
\mathcal{F}[u \mapsto \operatorname{sech}(u)](t) = \pi \operatorname{sech} \left(\frac{\pi}{2}t\right)
$$
Hence, by the Poisson summation formula, for every $x \in \mathbb{R}$ we get
\begin{align*}
\sum_{k \in \mathbb{Z}} \operatorname{sech}(xk) &=\ \sum_{k \in \mathbb{Z}} \mathcal{F}\left[y \mapsto \operatorname{sech}(xy)\right](\eta)\big|_{\eta = k}\\
&=\ \sum_{k\in \mathbb{Z}} \int_\mathbb{R} \operatorname{sech}(u) e^{-i(2\pi/x) \cdot \eta} \frac{\mathrm{d}u}{x}\big|_{\eta = k}\\
&=\ \frac{\pi}{|x|} \sum_{k \in \mathbb{Z}} \operatorname{sech}\left(\frac{\pi^2}{x}\eta\right) \big|_{\eta = k},
\end{align*}
which results in Equation (\ref{eq:functional_F}). Now, we substitute $x = \sqrt{r} \cdot \pi$. For every $r \in \mathbb{N}$, Ramanujan's formula \ref{form:ramanujan_1} yields
$$
\sum_{k\in \mathbb{Z}} \operatorname{sech}\left(\frac{\pi}{\sqrt{r}} k \right) = \sqrt{r} \left(\sum_{k \in \mathbb{Z}} \operatorname{sech}\left(\sqrt{r} \cdot \pi k \right)\right) = \frac{2\sqrt{r}}{\pi}K\left(k_r\right),
$$
as claimed.
\end{proof}

For the scaling factor $1/\sqrt{r}$, some relevant normalization constants can easily be calculated by multiplying the corresponding normalizing constants for the scaling factor $\sqrt{r}$ in Table \ref{tab:sqrt_r_sech} with the number $1/\sqrt{r}$. This yields some singular values and matching probability mass functions $\left(q_{r,k}\right)_{k \in \mathbb{Z}}$ as follows:

\renewcommand{\arraystretch}{2}
\begin{table}[t]
    \centering
    \begin{tabular}{c|c|r}
$q_{r,k}$  &   $ \lambda^*(r) = k_r$ & $\frac{\pi}{\sqrt{r} \cdot 2K(k_r)}\operatorname{sech} (\sqrt{r}\pi k)$\\
\hline
$\mathbf{q_{1,k}}$  & $\frac{1}{\sqrt{2}}$ & $\frac{1}{\sqrt{2}G} \operatorname{sech} \left(\pi k\right)\quad  $\\
$q_{2,k}$  & $\sqrt{2}-1$ & $\frac{4\sqrt{\pi^3}}{\sqrt{\sqrt{2}+1} \cdot \Gamma(1/8)\Gamma(3/8)} \operatorname{sech} \left(\frac{\pi}{\sqrt{2}} k\right)$\\
$q_{3,k}$  & $\frac{\sqrt{2}\left(\sqrt{3}-1\right)}{4}$ & $\frac{\sqrt[4]{2^3} \cdot \pi^2 }{\sqrt[4]{3^3} \cdot \Gamma^3(1/3)}\operatorname{sech} \left(\frac{\pi}{\sqrt{3}} k\right)$\\
$\mathbf{q_{4,k}}$  & $3-2\sqrt{2}$ & $\frac{1}{(1 + \sqrt{2})G}\operatorname{sech} \left(\frac{\pi}{2} k\right)\quad$\\
$q_{5,k}$  & $\frac{1}{\sqrt{2}} - \sqrt{\sqrt{5}-2}$ & $\frac{\sqrt{8 \pi^3}}{\sqrt[4]{2+\sqrt{5}} \sqrt{\Gamma(1/20)\Gamma(3/20)\Gamma(7/20)\Gamma(9/20)}} \operatorname{sech} \left(\frac{\pi}{\sqrt{5}} k\right)$\\
$q_{6,k}$  & $\left(2-\sqrt{3}\right)\left(\sqrt{3}-\sqrt{2}\right)$ & $\frac{\sqrt{16 \pi^3}}{\sqrt{(1 +\sqrt{2}+\sqrt{6}) \cdot \Gamma(1/24)\Gamma(5/24)\Gamma(7/24)\Gamma(11/24)}} \operatorname{sech} \left(\frac{\pi}{\sqrt{6}} k\right)$\\
$q_{7,k}$  & $\frac{3-\sqrt{7}}{8}$ & $\frac{2\pi^2}{\sqrt[4]{7} \cdot \Gamma(1/7)\Gamma(2/7)\Gamma(4/7)} \operatorname{sech} \left(\frac{\pi}{\sqrt{7}} k\right)$\\
% $q_{8,k}$  & $\frac{\pi}{2K\left(\left(1 + \sqrt{2} -\sqrt{2\sqrt{2}+2}\right)^2\right)} \operatorname{sech} (\sqrt{8} \pi k)$ & $\frac{2^{1/4}\cdot 8 \pi^{3/2}}{\left(\sqrt{2}+1\right)^{1/4} \sqrt{2\sqrt{2} + \sqrt{1 + 5 \sqrt{2}}} \cdot \Gamma(1/8)\Gamma(3/8)} \operatorname{sech} (\sqrt{8}\pi k)$\\
$\mathbf{q_{9,k}}$  & $\frac{\left(\sqrt{2}-\sqrt[4]{3}\right)\left(\sqrt{3}-1\right)}{2}$ & $\frac{1}{\left(\sqrt[4]{3^3} + \sqrt[4]{3} \right) G} \operatorname{sech} \left(\frac{\pi}{3} k\right)\quad$\\
$\mathbf{q_{16,k}}$  & $\left(1 + \sqrt{2}\right)^2 \left(\sqrt[4]{2} - 1\right)^4$ & $\frac{1}{\left(1 + \sqrt[4]{2} \right)^2 G} \operatorname{sech} \left(\frac{\pi}{4} k\right)\quad$\\
$\mathbf{q_{25,k}}$  & $\frac{\left(3 - 2 \sqrt[4]{5}\right)\left(\sqrt{5} - 2\right)}{\sqrt{2}}$ & $\frac{1}{\left(2\sqrt{2} + \sqrt{10} \right) G} \operatorname{sech} \left(\frac{\pi}{5} k\right) \quad$
\end{tabular}
    \caption{Note, that the probability mass functions highlighted in \textbf{bold} have simple normalization constants with a $1$ in the nominator and scalar times  Gau\ss' constant $G=\varpi/\pi$.}
    \label{tab:sqrt_r_sech}
\end{table}

\vspace{3cm}

\section{The First Discrete Bilateral Hyperbolic Secant Distribution}

We calculate some moments for the first discrete bilateral hyperbolic secant distribution.

\begin{theorem} \label{thm_Var_1_b}
Let $X\sim\mathsf{Sech}_{1,b}$ be a discrete random variable which is distributed according to the first discrete bilateral hyperbolic secant distribution. The odd moments vanish and we have
\begin{align*}    
(i) \quad \mathbb{E}\left[X^2 \right] =&\ \mathrm{Var}(X) = \left(\frac{G}{2}\right)^2 \approx 0.1741505,\\
(ii) \quad \mathbb{E}\left[X^4 \right] =&\ 9\left(\frac{G}{2}\right)^4 \approx 0.2729555,\\
(iii) \quad \mathbb{E}\left[X^6 \right] =&\ 153\left(\frac{G}{2}\right)^6 \approx 0.8081008,\\
(iv) \quad \mathbb{E}\left[X^8 \right] =&\ 4977\left(\frac{G}{2}\right)^8 \approx 4.5779016,
\end{align*}

with Gau\ss' constant $G = \varpi/\pi$ and the lemniscate constant $\varpi = \Gamma^2(1/4)/(2\sqrt{2\pi})$.
\end{theorem}

\begin{proof}
Let $C_{1,b} = 1/(\sqrt{2}G)$. We get $\mathbb{E}[X^n] = 0$ for odd $n$ due to symmetry. The other formulas directly follow by Ramanujan's \textit{Entries 17 (ii)-(v)} in Chapter 17 of his second notebook, proven by Berndt on pp. 138-139 in \cite{ramanujan_1991}: For $j\in \{2,4,6,8\}$, we write
$$
\mathbb{E}\left[X^{j}\right] = C_{1,b} \sum_{k\in\mathbb{Z}} k^{j} \operatorname{sech}(\pi k) = \frac{\sqrt{2}}{G} \sum_{n=1}^\infty n^{j} \operatorname{sech}(\pi n).
$$
We translate Ramanujan's expressions into the common notation via $x = k^2$, $y = \pi K/K'$, and $z=2 K/\pi$ and use the elliptic lambda function $\lambda^\star(k_1) = 1/\sqrt{2}$ with $K/K'=1$ and $K(k_1)=\varpi/\sqrt{2}$. This yields $x=1/2$, $y=\pi$ and $z=\sqrt{2}G$. We have

$$
\sum_{n=1}^\infty n^{2} \operatorname{sech}(y n) = \frac{1}{8} z^3 x
\ \Rightarrow \
\mathbb{E}\left[X^2\right] = \frac{1}{4} G^2,
$$
$$
\sum_{n=1}^\infty n^{4} \operatorname{sech}(y n) = \frac{1}{2} z^5 \left(\frac{x}{4} + \left(\frac{x}{4} \right)^2 \right)
\ \Rightarrow \
\mathbb{E}\left[X^4\right] = \frac{9}{16} G^4,
$$
$$
\sum_{n=1}^\infty n^{6} \operatorname{sech}(y n) = \frac{1}{2} z^7 \left(\frac{x}{4} + 11\left(\frac{x}{4}\right)^2 + \left(\frac{x}{4} \right)^3 \right)
\ \Rightarrow \
\mathbb{E}\left[X^6\right] = \frac{153}{64} G^6,
$$
$$
\sum_{n=1}^\infty n^{8} \operatorname{sech}(y n) = \frac{1}{2} z^9 \left(\frac{x}{4} + 57\left(\frac{x}{4}\right) + 102 \left(\frac{x}{4} \right)^3 + \left(\frac{x}{4} \right)^4 \right)
\ \Rightarrow \
\mathbb{E}\left[X^8\right] = \frac{4977}{256} G^8,
$$
\vspace{0.7cm}
as claimed.
\end{proof}

Now, we can deduce some exact expressions for series of odd powers of the hyperbolic secant via the Poisson summation.

\begin{theorem}
The following identities hold:

\begin{align*}
(i) \quad \sum_{k\in\mathbb{Z}} \mathrm{sech}^3 (\pi k) =&\ \frac{\sqrt{2}}{2}\left(G^3 + G\right) \approx 1.0012841,\\
(ii) \quad \sum_{k \in \mathbb{Z}} \mathrm{sech}^5(\pi k) =&\ \frac{\sqrt{2}}{24}(9 G^5 + 10 G^3 + 9 G) \approx 1.0000096,\\
(iii) \quad \sum_{k \in \mathbb{Z}} \mathrm{sech}^7(\pi k) =&\ \frac{\sqrt{2}}{720}\left(153 G^7 + 315 G^5 + 259 G^3 + 225 G\right) \approx 1.0000000711,\\
(iv) \quad \sum_{k \in \mathbb{Z}} \mathrm{sech}^9(\pi k) =&\ \frac{\sqrt{2}}{42320}\left(4977 G^9 + 12852 G^7 + 17766 G^5 + 12916 G^3 + 11025 G\right)\\
\approx&\ 1.0000000005,
\end{align*}
with Gau\ss' constant $G = \varpi/\pi$ and the lemniscate constant $\varpi = \Gamma^2(1/4)/(2\sqrt{2\pi})$.
\end{theorem}

\begin{proof}
Let $C_{1,b} = 1/(\sqrt{2}G)$. For $n \in \{2, 4, 6, 8\}$, we apply Poisson's summation formula to the identities $C_{1,b} \sum_{k\in\mathbb{Z}} k^{n} \mathrm{sech}(\pi k)$, i.\,e. the $n$-th moments of a discrete random variable which is distributed according to the first bilateral hyperbolic secant distribution. Define the function $s(x) = x^{n} \cdot \mathrm{sech}(\pi x)$ and the Fourier transform
\begin{align*}
\mathcal{F}[s](\xi) =&\ \int_{-\infty}^\infty s(x) e^{-2\pi i x \cdot \xi} \mathrm{d}x.
\end{align*}
It is easy to see that for every $n\in\mathbb{N}$, we get
$$
\mathcal{F}[x^n f(x)](\xi) = \frac{1}{(-2\pi i)^n} \frac{\partial^n}{\partial\xi^n} [\mathcal{F}[f(x)](\xi)]
$$
for any Schwartz function $f \in \mathcal{S}(\mathbb{R})$ and it is well-known that
$$
\mathcal{F}[\mathrm{sech}(\pi x)](\xi) = \mathrm{sech(\pi \xi)}
$$
holds, i.\,e. this version of the hyperbolic secant is a fix-point of this particular form of the Fourier transform. Now, Poisson's summation formula states
$$\sum_{k\in \mathbb{Z}}s(k) = \sum_{\kappa \in \mathbb{Z}}\mathcal{F}[s](\kappa).$$
\textit{(i)} For $n=2$, we have 
\begin{align*}
\mathcal{F}[s](\xi) = -\frac{1}{4 \pi ^2} \frac{\partial^2}{\partial\xi^2} [\mathrm{sech}(\pi \xi)]
= \frac{1}{4}(2 \cdot \mathrm{sech}^3(\pi \xi) - \mathrm{sech}(\pi \xi)). 
\end{align*}
%\cdot 2^{1-2r} \left(\sum_{k=1}^{r/2} (-1)^k T(r+1,k) \cosh (2k\pi \xi) + \frac{1}{2} T\left(r+1, \frac{r}{2}+1\right) \right) \mathrm{sech}^{r+1}(\pi \xi), with Euler's numbers $T(n,k)$ of type B, cf. https://oeis.org/A060187.
This results in
$$
\mathbb{E}\left[X^2\right] = \frac{C_{1,b}}{2} \sum_{\kappa \in \mathbb{Z}} \mathrm{sech}^3(\pi \kappa) - \frac{C_{1,b}}{4} \sum_{\kappa \in \mathbb{Z}} \mathrm{sech}(\pi \kappa).
$$
But the last summand can be explicitly calculated thanks to one of Ramanujan's identities, cf. p. 138 of \cite{ramanujan_1991}. It was already used in the proof of Theorem \ref{thm:pmfs}. One has
$$
\sum_{\kappa \in \mathbb{Z}} \mathrm{sech}(\pi \kappa) = 2 \sum_{n=1}^\infty \mathrm{sech}(\pi n) + 1 = 2 \frac{\sqrt{2}G - 1}{2}  + 1 =\sqrt{2}G = \frac{1}{C_{1,b}}.
$$
This finally leads to
$$
\mathrm{Var}(X) = \mathbb{E}\left[X^2\right] = \frac{1}{2\sqrt{2}G} \sum_{\kappa \in \mathbb{Z}} \mathrm{sech}^3(\pi \kappa)  - \frac{1}{4}.
$$
But, result \textit{(i)} from Theorem \ref{thm_Var_1_b} implies 
$$
\mathrm{Var}(X) = \frac{1}{4}G^2,
$$
which holds if and only if
$$
\sum_{k \in \mathbb{Z}} \mathrm{sech}^3(\pi k) = \frac{G^3 + G}{\sqrt{2}}.
$$
\textit{(ii)} For $n=4$, we have 
\begin{align*}
\mathcal{F}[s](\xi) =&\ \frac{1}{16 \pi ^4} \frac{\partial^4}{\partial\xi^4} \left[ \mathrm{sech}(\pi \xi)\right] = \frac{3}{2} \operatorname{sech}^5(\pi \xi) - \frac{5}{4}\operatorname{sech}^3 (\pi \xi) + \frac{1}{16} \operatorname{sech}(\pi \xi).
\end{align*}
%\cdot 2^{1-2r} \left(\sum_{k=1}^{r/2} (-1)^k T(r+1,k) \cosh (2k\pi \xi) + \frac{1}{2} T\left(r+1, \frac{r}{2}+1\right) \right) \mathrm{sech}^{r+1}(\pi \xi), with Euler's numbers $T(n,k)$ of type B, cf. https://oeis.org/A060187.
The Poisson summation results in
$$
\mathbb{E}\left[X^4\right] = \frac{3}{2} C_{1,b} \sum_{\kappa \in \mathbb{Z}} \mathrm{sech}^5(\pi \kappa) - \frac{5}{4} C_{1,b} \sum_{\kappa \in \mathbb{Z}} \mathrm{sech}^3(\pi \kappa) + \frac{1}{16} C_{1,b} \sum_{\kappa \in \mathbb{Z}} \mathrm{sech}(\pi \kappa).
$$
But the last two series can be explicitly calculated by part \textit{(i)}, viz.
$$
\sum_{\kappa \in \mathbb{Z}} \mathrm{sech}^3(\pi \kappa) = \frac{G^3 + G}{\sqrt{2}}, \quad \sum_{\kappa \in \mathbb{Z}} \mathrm{sech}(\pi \kappa) = \frac{1}{C_{1,b}} = \sqrt{2}G.
$$
This yields
$$
\mathbb{E}\left[X^4\right] = \frac{3}{2\sqrt{2}G} \sum_{\kappa \in \mathbb{Z}} \mathrm{sech}^5(\pi \kappa) - \frac{5}{4\sqrt{2}G} \frac{G^3 + G}{\sqrt{2}} + \frac{1}{16}.
$$
But, result \textit{(ii)} from Theorem \ref{thm_Var_1_b} tells us
$$
\mathbb{E}\left[X^4\right] = \frac{9}{16}G^4,
$$
which holds if and only if
$$
\sum_{k \in \mathbb{Z}} \mathrm{sech}^5(\pi k) = \frac{\sqrt{2}}{24}(9 G^5 + 10G^3+ 9G).
$$

\textit{(iii)} For $n=6$, we have
\begin{align*}
\mathcal{F}[s](\xi) =&\ -\frac{1}{64 \pi ^6} \frac{\partial^6}{\partial\xi^6} \left[ \mathrm{sech}(\pi \xi)\right]\\
=&\ \frac{45}{4} \operatorname{sech}^7(\pi \xi) - \frac{105}{8}\operatorname{sech}^5 (\pi \xi) + \frac{91}{32} \operatorname{sech}^3(\pi \xi) - \frac{1}{64} \operatorname{sech}(\pi \xi).
\end{align*}
%\cdot 2^{1-2r} \left(\sum_{k=1}^{r/2} (-1)^k T(r+1,k) \cosh (2k\pi \xi) + \frac{1}{2} T\left(r+1, \frac{r}{2}+1\right) \right) \mathrm{sech}^{r+1}(\pi \xi), with Euler's numbers $T(n,k)$ of type B, cf. https://oeis.org/A060187.
The Poisson summation results in
\begin{align*}
\mathbb{E}\left[X^6\right] =\ &\frac{45}{4} C_{1,b} \sum_{\kappa \in \mathbb{Z}} \operatorname{sech}^7(\pi \kappa) - \frac{105}{8} C_{1,b} \sum_{\kappa \in \mathbb{Z}} \operatorname{sech}^5 (\pi \kappa)\\
&+ \frac{91}{32} C_{1,b} \sum_{\kappa \in \mathbb{Z}} \operatorname{sech}^3(\pi \kappa) - \frac{1}{64} C_{1,b} \sum_{\kappa \in \mathbb{Z}} \operatorname{sech}(\pi \kappa).
\end{align*}
But the last four series can be explicitly calculated by the parts \textit{(i)} and \textit{(ii)}, viz.
\begin{align*}
& \sum_{k \in \mathbb{Z}} \mathrm{sech}^5(\pi k) = \frac{\sqrt{2}}{24}\left(9 G^5 + 10 G^3 + 9 G\right), \\ & \sum_{\kappa \in \mathbb{Z}} \mathrm{sech}^3(\pi \kappa) = \frac{\sqrt{2}}{2} \left(G^3+G\right), \\ & \sum_{\kappa \in \mathbb{Z}} \mathrm{sech}(\pi \kappa) = \sqrt{2}G.
\end{align*}
This yields
$$
\mathbb{E}\left[X^6\right] = \frac{45}{4\sqrt{2}G} \sum_{\kappa \in \mathbb{Z}} \mathrm{sech}^7(\pi \kappa) - \frac{105}{192} \left(9 G^4 + 10 G^2 + 9\right) + \frac{91}{64} \left(G^2 +1 \right) - \frac{1}{64}.
$$
But, result \textit{(ii)} from Theorem \ref{thm_Var_1_b} tells us
$$
\mathbb{E}\left[X^6\right] = \frac{153}{64}G^6,
$$
which holds if and only if
$$
\sum_{k \in \mathbb{Z}} \mathrm{sech}^7(\pi k) = \frac{\sqrt{2}}{720}\left(153 G^7 + 315 G^5 + 259 G^3 + 225 G\right).
$$

\textit{(iv)} For $n=8$, we have
\begin{align*}
\mathcal{F}[s](\xi) =&\ -\frac{1}{256 \pi^8} \frac{\partial^8}{\partial\xi^8} \left[ \mathrm{sech}(\pi \xi)\right]\\
=&\ \frac{315}{2} \operatorname{sech}^9(\pi \xi) - \frac{945}{4} \operatorname{sech}^7(\pi \xi) + \frac{1449}{16}\operatorname{sech}^5 (\pi \xi) - \frac{205}{32} \operatorname{sech}^3(\pi \xi) + \frac{1}{256} \operatorname{sech}(\pi \xi).
\end{align*}
%\cdot 2^{1-2r} \left(\sum_{k=1}^{r/2} (-1)^k T(r+1,k) \cosh (2k\pi \xi) + \frac{1}{2} T\left(r+1, \frac{r}{2}+1\right) \right) \mathrm{sech}^{r+1}(\pi \xi), with Euler's numbers $T(n,k)$ of type B, cf. https://oeis.org/A060187.
The Poisson summation formula results in
\begin{align*}
\mathbb{E}\left[X^8\right] = \sum_{\kappa \in \mathbb{Z}} \Bigg[ & \frac{315}{2\sqrt{2}G}\operatorname{sech}^9(\pi \kappa) - \frac{945}{4\sqrt{2}G} \operatorname{sech}^7 (\pi \kappa) + \frac{1449}{16\sqrt{2}G} \operatorname{sech}^5(\pi \kappa)\\ 
&- \frac{205}{32\sqrt{2}G} \operatorname{sech}^3(\pi \kappa) + \frac{1}{256 \sqrt{2}G} \operatorname{sech}(\pi \kappa) \Bigg].
\end{align*}
But the last four series can be explicitly calculated by our parts \textit{(i)}, \textit{(ii)} and \textit{(iii)}, viz.
\begin{align*}
& \sum_{k \in \mathbb{Z}} \mathrm{sech}^7(\pi k) = \frac{\sqrt{2}}{720}\left(153 G^7 + 315 G^5 + 259 G^3 + 225 G\right)\\
& \sum_{k \in \mathbb{Z}} \mathrm{sech}^5(\pi k) = \frac{\sqrt{2}}{24}\left(9 G^5 + 10 G^3 + 9 G\right), \\ & \sum_{\kappa \in \mathbb{Z}} \mathrm{sech}^3(\pi \kappa) = \frac{\sqrt{2}}{2} \left(G^3+G\right), \\ & \sum_{\kappa \in \mathbb{Z}} \mathrm{sech}(\pi \kappa) = \sqrt{2}G.
\end{align*}

But, result \textit{(ii)} from Theorem \ref{thm_Var_1_b} tells us
$$
\mathbb{E}\left[X^8\right] = \frac{4977}{256}G^8,
$$
which holds if and only if
$$
\sum_{k \in \mathbb{Z}} \mathrm{sech}^9(\pi k) = \frac{\sqrt{2}}{42320}\left(4977 G^9 + 12852 G^7 + 17766 G^5 + 12916 G^3 + 11025 G\right),
$$
as claimed.
\end{proof}

\vspace{1cm}

\section{The Second Discrete Bilateral Hyperbolic Secant Distribution}

Ramanujan's {Entries 16 (x) - (xiii) in Chapter 17 of his second notebook, proven by Berndt on pp. 137-138 in \cite{ramanujan_1991}, can be statistically interpreted as follows: If you randomly generate odd integers with even multiplicity by using the second bilateral hyperbolic secant distribution, you can calculate the expected outcome up to the eighth order:

\begin{theorem} 
Let $Y\sim\mathsf{Sech}_{2,b}$ be a discrete random variable which is distributed according to the second bilateral hyperbolic secant distribution. We have
\begin{align*}    
(i) \quad \mathbb{E}\left[(2Y+1)^2 \right] =&\ 2G^2 \approx 1.3932039,\\
(ii) \quad \mathbb{E}\left[(2Y+1)^4 \right] =&\ 12 G^4 \approx 5.8230516,\\
(iii) \quad \mathbb{E}\left[(2Y+1)^6 \right] =&\ 216 G^6 \approx 73.0142850,\\
(iv) \quad \mathbb{E}\left[(2Y+1)^8 \right] =&\ 7056 G^8 \approx 1661.4885492,
\end{align*}

with Gau\ss' constant $G = \varpi/\pi$ and the lemniscate constant $\varpi = \Gamma^2(1/4)/(2\sqrt{2\pi})$.
\end{theorem}

\begin{proof}
Let $C_{2,b} = 1/G$. We get $\mathbb{E}[Y^j] = 0$ for odd $j \in \mathbb{N}$ due to symmetry. The other formulas directly follow by Ramanujan's \textit{Entries 16 (x)-(xiii)} in Chapter 17 of his second notebook, proven by Berndt on pp. 137-138 in \cite{ramanujan_1991}: For $j\in \{2,4,6,8\}$, we write
$$
\mathbb{E}\left[Y^{j}\right] = C_{2,b} \sum_{k\in\mathbb{Z}} (2k+1)^{j} \operatorname{sech}\left(\frac{\pi}{2} (2k+1)\right) = \frac{2}{G} \sum_{n=0}^\infty (2n+1)^{j} \operatorname{sech}\left(\frac{\pi}{2} (2n+1)\right).
$$
We translate Ramanujan's expressions into the common notation via $x = k^2$, $y = \pi K/K'$, and $z=2 K/\pi$ and use the elliptic lambda function $\lambda^\star(k_1) = 1/\sqrt{2}$ with $K/K'=1$ and $K(k_1)=\varpi/\sqrt{2}$. This yields $x=1/2$, $y=\pi$ and $z=\sqrt{2}G$. We have

\begin{align*}
\textit{(i)} \quad
\sum_{n=0}^\infty (2n+1)^{2} \operatorname{sech}\left(\frac{1}{2} (2n+1)y\right) &= \frac{1}{2} z^3 \sqrt{x}\\
& \Rightarrow \mathbb{E}\left[(2Y+1)^2 \right]= 2 G^2,\\
\textit{(ii)} \quad
\sum_{n=0}^\infty (2n+1)^{4} \operatorname{sech}\left(\frac{1}{2} (2n+1)y\right) &= \frac{1}{2} z^5 (1+4x) \sqrt{x}\\
& \Rightarrow \mathbb{E}\left[(2Y+1)^4 \right] = 12 G^4,\\
\textit{(iii)} \quad
\sum_{n=0}^\infty (2n+1)^{6} \operatorname{sech}\left(\frac{1}{2} (2n+1)y\right) &= \frac{1}{2} z^7 (1+11(4x) + (4x)^2) \sqrt{x}\\
& \Rightarrow \mathbb{E}\left[(2Y+1)^6 \right] = 216 G^6,\\
\textit{(iv)} \quad
\sum_{n=0}^\infty (2n+1)^{8} \operatorname{sech}\left(\frac{1}{2} (2n+1)y\right) &= \frac{1}{2} z^9 (1 + 102(4x) + 57(4x)^2 + (4x)^3) \sqrt{x}\\
& \Rightarrow \mathbb{E}\left[(2Y+1)^8 \right] = 7056 G^8,
\end{align*}
as claimed.
\end{proof}

\section{Application to Contour Integrals}

We apply our result to complex analysis and calculate some contour integral on the critical line in terms of the variance of our first discrete bilateral hyperbolic secant distribution. We integrate the product of the analytic continuation of Dirichlet's beta function $\beta$ and the meromorphic continuations of Euler's Gamma function $\Gamma$ and Riemann's zeta function $\zeta$. Further, we use Euler's numbers which are defined via
$$
\operatorname{sech}(t) = \sum_{n=0}^\infty \frac{E_n}{n!} \cdot t^n,$$
vanishing for odd $n$.

\begin{theorem}
Let $\Xi(s):= \pi^{-(s+2)} \cdot \beta(s+2) \cdot \Gamma(s+2) \cdot \zeta(s)$. On the critical line, one has
\begin{align*}
\int_{1/2-i\cdot \infty}^{1/2 + i \cdot \infty} \Xi(s)\, \mathrm{d}s = i \left( \frac{\varpi}{\sqrt{2}} \cdot \mathbb{E}\left[X^2 \right] - \frac{\pi}{8} \right) = i \left( \frac{\varpi}{\sqrt{2}} \left(\frac{G}{2}\right)^2 - \frac{\pi}{8} \right) \approx&\ -i \cdot 0.0698111,
\end{align*}
with Gau\ss' constant $G = \varpi/\pi$ and the lemniscate constant $\varpi = \Gamma^2(1/4)/(2\sqrt{2\pi})$.
\end{theorem}

\begin{proof}
Let $X\sim\mathsf{Sech}_{1,b}$ be a discrete random variable which is distributed according to the first bilateral hyperbolic secant distribution. We expand the second moment of $X$. Let $C_{1,b} = 1/(\sqrt{2}G)$. By definition, one has
$$
\mathbb{E}\left[X^2\right] = C_{1,b} \sum_{k \in \mathbb{Z}} k^2\cdot \operatorname{sech}\left(\pi k\right) = 2C_{1,b} \sum_{n=1}^\infty n^2 \cdot \operatorname{sech}(\pi n).
$$
For $x > 0$, one has $|\exp(-2x)| < 1$. Hence, the geometric series yields the identities
$$
\operatorname{sech}(x) = \frac{2e^{-x}}{1+e^{-2x}} = 2e^{-x} \sum_{k=0}^\infty (-1)^k e^{-2kx} = 2 \sum_{k=0}^\infty (-1)^k e^{-(2k+1)x}.
$$
Now, the gamma function motivates the introduction of the Mellin transform $\mathcal{M}$ on a strip $\mathbb{S} := \{0 < \Re(\sigma) < b\}$. For every $f \in L^1_{\mathrm{loc}}((\mathbb{S},\infty), \mathbb{C})$, it is defined as
$$
\mathcal{M}\left\{f\right\}(s) = \int_0^\infty x^{s-1} f(x) \, \mathrm{d}x, \quad \mathrm{with} \quad \mathcal{M}\left\{\left[x \mapsto e^{-x}\right]\right\} = \int_0^\infty x^{s-1}e^{-x} \, \mathrm{d}x = \Gamma(s).
$$
The gamma function is analytic on $\mathbb{S}$ and Stirling's formula
\begin{equation}\label{Stirlings_formula}
|\Gamma(\sigma + i \cdot t)| \in O\left(\sqrt{2\pi} \cdot |t|^{\sigma-1/2} e^{-|t| \pi/2}\right) (|t| \rightarrow \infty)
\end{equation}
holds uniformly for any $\sigma > 0$. Hence, the Mellin inversion theorem yields the existence of the Cahen-Mellin integral
$$
e^{-x} = \mathcal{M}^{-1}\left\{\Gamma\right\}(x) = \frac{1}{2\pi i} \int_{c-i \cdot\infty}^{c+i \cdot\infty} \Gamma(s) \cdot x^{-s} \, \mathrm{d}s.
$$
For any fixed $c>0$ we conclude
$$
\mathbb{E}\left[X^2\right] = 4C_{1,b} \sum_{n=1}^\infty n^2 \sum_{k=0}^\infty (-1)^k \frac{1}{2\pi i} \int_{c-i \cdot\infty}^{c+i \cdot\infty} \frac{\Gamma(s)}{(2k+1)^{s}} \cdot (\pi n)^{-s}\, \mathrm{d}s.
$$
Next, we want to interchange the limits. On the line $\Re(s)=c>3$, we get
$$
\left|(-1)^k \cdot \frac{\Gamma(s)\cdot n^{2-s}}{(\pi (2k+1))^{s}}\right| = \frac{\Gamma(c)\cdot n^{2-c}}{\left(\pi (2k+1)\right)^{c}}
$$
and by recalling Stirling's formula (\ref{Stirlings_formula}) we infer
$$
\int_{c-i \cdot\infty}^{c+i \cdot\infty} |\Gamma(s)|\, \mathrm{d}s < \infty, \quad \sum_{k=0}^\infty (2k+1)^{-c} < \infty, \quad \sum_{n=1}^\infty n^{2-c} < \infty.
$$
Thus, Tonelli's theorem results in
\begin{align*}
\mathbb{E}\left[X^2\right] =&\ \frac{4C_{1,b}}{2\pi i} \int_{c-i \cdot\infty}^{c+i \cdot\infty} \frac{\Gamma(s)}{\pi^s} \sum_{k=0}^\infty \frac{(-1)^k}{(2k+1)^{s}} \sum_{n=1}^\infty \frac{1}{n^{s-2}}\, \mathrm{d}s \\
=&\ \frac{4C_{1,b}}{2\pi i} \int_{c-i \cdot\infty}^{c+i \cdot\infty} \underbrace{\pi^{-s} \cdot \beta(s) \cdot \Gamma(s) \cdot \zeta(s - 2)}_{=:\, \Phi(s)} \, \mathrm{d}s.
\end{align*}

We want to solve the integral by the method of residues by shifting the vertical line of integration to the left. Therefore we consider the analytic continuation of $\beta$ and the meromorphic continuations of $\Gamma$, and $\zeta$. In the representation $\Phi(s)$ of the integrand, the simple poles of $\Gamma$ and the trivial zeros of $\zeta(\cdot  - 2)$ cancel each other out for even non-positive integers $s\in \{2-2n\}_{n\in \mathbb{N}}$. In addition, the simple poles of $\Gamma$ and the simple zeros of $\beta$ cancel each other out for odd negative integers $s \in \{1-2n\}_{n \in \mathbb{N}}$. Hence, all $\Gamma$-poles in $s\in\{-n\}_{n\in \mathbb{N}_0}$ are canceled out and only the simple pole at $s = \varrho = 3$ from $\zeta(\cdot - 2)$ remains with residue
$$
\operatorname{Res}_{\Phi}(\varrho = 3) = \pi^{-3} \cdot \beta(3) \cdot \Gamma(3) \cdot 1 = \frac{2!}{\pi^{3}} \beta(3) = \frac{- E_2}{4^2} = \frac{1}{4^2},
$$
since the identity
$$
\beta(r+1)=\frac{(-1)^{\frac{r}{2}} E_{r}}{2\cdot r!}\left(\frac\pi2\right)^{r+1}
$$
with Euler's numbers $E_r$ holds for even $r \in \mathbb{N}$ and one has $E_2 = -1$.\\

We introduce a new contour for the integral of $\Phi$. Let $c=3+\varepsilon$ for some $\varepsilon>0$ and let $R= R_{5/2,c}(\pm T)$ be the rectangle with vertices $c \pm i\cdot T$ and $5/2 \pm i \cdot T$. Then,by the residue theorem, one has
$$
\frac{4 C_{1,b}}{2\pi i}\int_{R} \Phi(s) \, \mathrm{d}s = 4 C_{1,b} \sum_{\varrho \in R} \operatorname{Res}_\Phi(\varrho) = \frac{C_{1,b}}{4},
$$
while
\begin{align*}
&\ \frac{4 C_{1,b}}{2\pi i}\int_{c-i\cdot T}^{c+i\cdot T} \Phi(s) \mathrm{d}s\\
=&\ \frac{4 C_{1,b}}{2\pi i} \left[\underbrace{\int_{5/2 + i\cdot T}^{c + i\cdot T} \Phi(s) \mathrm{d}s}_{=: I_1} + \underbrace{\int_{5/2-i\cdot T}^{5/2+i\cdot T} \Phi(s) \mathrm{d}s}_{=: I_2} + \underbrace{\int_{5/2 - i \cdot T}^{c- i \cdot T} \Phi(s) \mathrm{d}s}_{=: I_3}\right] + \frac{C_{1,b}}{4}.
\end{align*}

We show that the integrals over the horizontal edges vanish: For $\sigma > 1$, the beta function converges absolutely. Hence, one has
$$
\lvert I_1\rvert, \lvert I_3\rvert \le \int_{5/2 \pm i\cdot T}^{c \pm i\cdot T} \lvert \pi^{-s} \beta(s) \Gamma(s) \zeta(s-2) \rvert \, \mathrm{d}s \lesssim \int_{5/2 \pm i\cdot T}^{c \pm i\cdot T} \lvert \Gamma(s) \rvert \cdot \zeta(\sigma-2) \, \mathrm{d}s.
$$
The zeta term grows at most exponentially and Stirling's formula (\ref{Stirlings_formula}) results in 
$$\lvert \Phi(\sigma + i \cdot T)\rvert \in O\left(T^{\sigma/2} e^{-T\pi/2}\right)(T \rightarrow \infty).$$
Therefore, $\sup_{\sigma \in [5/2, 3+\varepsilon]}\lvert \Phi(\sigma + i \cdot T) \rvert \rightarrow 0$ as $T \rightarrow \infty$ and $I_1$ and $I_3$ uniformly vanish in the limit. Now, we consider the shifted function 
$$\Xi(s):= \Phi(s+2) = \pi^{-(s+2)} \cdot \beta(s+2) \cdot \Gamma(s+2) \cdot \zeta(s) $$ 
on the line $\Re(s)=1/2.$ First, we show that $\lim_{T \rightarrow \infty} I_2$ exists. The beta function converges absolutely for $\sigma > 1$. Further, Bourgain showed in \cite{bourgain_2016}, that
$$
\zeta(1/2 + i \cdot t) \in O\left(\lvert t\rvert^{\frac{13}{84}+\delta}\right) (\lvert t\rvert \rightarrow \infty)
$$
for any $\delta > 0$. Together with Stirling's formula, those facts result in
$$\Xi\left(\frac{1}{2} \pm i \cdot T\right) \in O\left(T^{2 + \frac{13}{84} + \delta} e^{-T \pi/2}\right) (T\rightarrow\infty).$$
Hence, due to the exponential decay, the expression 
$$\mathbb{E}\left[X^2 \right] = \frac{2 C_{1,b}}{\pi i} \int_{1/2-i\cdot \infty}^{1/2 + i \cdot \infty} \Xi(s) \mathrm{d}s + \frac{C_{1,b}}{4}$$
is well-defined. Now, Theorem \ref{thm_Var_1_b} says $\mathbb{E}\left[X^2\right] = (G/2)^2$. Hence, in terms of the lemniscate constant $\varpi$ and the lemniscate modulus $k=1/\sqrt{2}$, we get the formula
\begin{align*}
\int_{1/2-i\cdot \infty}^{1/2 + i \cdot \infty} \Xi(s) \mathrm{d}s = i \left( \frac{\varpi}{\sqrt{2}} \cdot \mathbb{E}\left[X^2 \right] - \frac{\pi}{8} \right)= i \left( \frac{\varpi}{\sqrt{2}} \left(\frac{G}{2}\right)^2 - \frac{\pi}{8} \right),
\end{align*}
as claimed.
\end{proof}


\begin{thebibliography}{9}

\bibitem[Bourgain, 2016]{bourgain_2016}
Bourgain, Jean L. (2016).
\textit{Decoupling, exponential sums and the Riemann zeta function.} 
Via: https://arxiv.org/pdf/1408.5794

%\bibitem{bruijn_58}
%De Bruijn, N. G. (1958). \textit{Asymptotic Methods in Analysis}. Dover Publications. ISBN 978-0-486-64221-5


\bibitem[Campbell, 2023]{campbell_2023}
Campbell, John M. (2023). 
\textit{Hyperbolic summations derived using the Jacobi functions dc and nc.}
Via: https://arxiv.org/pdf/2301.03738v1.

\bibitem[Kirschenhofer, Prodinger and Szpankowski, 1996]{kirschenhofer_1996}
\ Kirschenhofer, P.; Prodinger, H.; Szpankowski, W. (1996). \textit{Multidimensional Digital Searching and Some New Parameters in Tries}. International J. Foundation of of Computer Science, 4, 69--84.

\bibitem[Ramanujan and Berndt, 1991]{ramanujan_1991}
Ramanujan, Srinivasa; Berndt, Bruce (ed.) (1991).
\textit{Ramanujan's Notebooks. Part III.}
Springer.

\bibitem[Whittaker and Watson, 1920]{whittaker_1920}
Whittaker, Edmund T.; Watson, George N. (1920).
\textit{A Course of Modern Analysis; an Introduction to the General Theory of Infinite Processes and of Analytic Functions; with an Account of the Principal Transcendental Functions (3rd Edition).}
Cambridge University Press.
Via: https://archive.org/details/courseofmodernan00whit/page/n7/mode/2up

\end{thebibliography}

% $$\E[X^r] = C_{1,b} \cdot 2^{1-2r} \sum_{l = 1}^\infty \left(\sum_{k=1}^{r/2} (-1)^k T(r+1,k) \cosh (2k\pi l) + \frac{1}{2} T\left(r+1, \frac{r}{2}+1\right) \right) \mathrm{sech}^{r+1}(\pi l).$$

% \begin{theorem}
%     Let $X\sim\mathsf{Sech}_{1,b}$ be a discrete random variable which is distributed according to the \textit{bilateral hyperbolic secant distribution of the first kind}. Let $r\in\N$ be a fixed positive integer. Then the moments of order $r$ of $X$ can be represented by
%     \begin{align*}
%         &\mathbb{E}[X^r]\\
%          =&\ C_{1,b} \cdot 2^{1-2r} \sum_{l = 1}^\infty \left(\sum_{k=1}^{r/2} (-1)^k T(r+1,k) \cosh (2k\pi l) + \frac{1}{2} T\left(r+1, \frac{r}{2}+1\right) \right) \mathrm{sech}^{r+1}(\pi l),
%     \end{align*}
% with Euler's numbers $T(n,k)$ of type B, cf. https://oeis.org/A060187.
% \end{theorem}

% \begin{proof}
% \textbf{Step 1:} Let $r$ be a fixed positive integer. We calculate the $r$-th moments of a random variable $X \sim \mathsf{Sech}_{1,b}$. For odd $r$ one has $\mathbb{E}[X^r] = 0$ due to symmetry. So, let $r$ be even. By definition, one has
% $$
% \mathbb{E}[X^r] = C_{1,b} \sum_{k \in \mathbb{Z}}^\infty k^r\cdot \operatorname{sech}\left(k \pi\right) = 2C_{1,b} \sum_{n=1}^\infty n^r \cdot \operatorname{sech}(n\pi).
% $$
% For $x > 0$, one has $|\exp(-2x)| < 1$. Hence, the geometric series yields the identities
% $$
% \operatorname{sech}(x) = \frac{2e^{-x}}{1+e^{-2x}} = 2e^{-x} \sum_{k=0}^\infty (-1)^k e^{-2kx} = 2 \sum_{k=0}^\infty (-1)^k e^{-(2k+1)x}.
% $$
% \textbf{Step 2:} Now, the gamma function motivates the introduction of the \textit{Mellin transform} $\mathcal{M}$ on a strip $\mathbb{S} := \{0 < \Re(\sigma) < b\}$. For every $f \in L^1_{\mathrm{loc}}((\mathbb{S},\infty), \mathbb{C})$, it is defined as
% $$
% \mathcal{M}\left\{f\right\}(s) = \int_0^\infty x^{s-1} f(x) \, \mathrm{d}x, \quad \mathrm{with} \quad \mathcal{M}\left\{\left[x \mapsto e^{-x}\right]\right\} = \int_0^\infty x^{s-1}e^{-x} \, \mathrm{d}x = \Gamma(s).
% $$
% The gamma function is analytic on $\mathbb{S}$ and Stirling's formula
% \begin{equation}\label{Stirlings_formula}
% |\Gamma(\sigma + i \cdot t)| \in O\left(\sqrt{2\pi} \cdot |t|^{\sigma-1/2} e^{-|t|\cdot \pi/2}\right) \cdot (|t| \rightarrow \infty)
% \end{equation}
% holds uniformly for any $\sigma > 0$. Hence, the Mellin inversion theorem yields the existence of the \textit{Cahen-Mellin integral}
% $$
% e^{-x} = \mathcal{M}^{-1}\left\{\Gamma\right\}(x) = \frac{1}{2\pi i} \int_{c-i \cdot\infty}^{c+i \cdot\infty} \Gamma(s) \cdot x^{-s} \, \mathrm{d}s.
% $$
% For any fixed $c>0$ we conclude
% $$
% \mathbb{E}\left[X^r\right] = 4C_{1,b} \sum_{n=1}^\infty n^r \sum_{k=0}^\infty (-1)^k \frac{1}{2\pi i} \int_{c-i \cdot\infty}^{c+i \cdot\infty} \frac{\Gamma(s)}{(2k+1)^{s}} \cdot (n\pi)^{-s}\, \mathrm{d}s.
% $$
% \textbf{Step 3:} Now, we want to interchange the limits. On the line $\Re(s)=c>r+1$ we get
% $$
% \left|(-1)^k \cdot \frac{\Gamma(s)}{((2k+1)\pi)^{s}} \cdot n^{r-s}\right| = \frac{\Gamma(c)}{\left((2k+1)\pi \right)^{c}} \cdot n^{r-c}
% $$
% and by recalling Stirling's formula (\ref{Stirlings_formula}) we infer
% $$
% \int_{c-i \cdot\infty}^{c+i \cdot\infty} |\Gamma(s)|\, \mathrm{d}s < \infty, \quad \sum_{k=0}^\infty (2k+1)^{-c} < \infty, \quad \sum_{n=1}^\infty n^{r-c} < \infty.
% $$
% Thus, Tonelli's theorem results in
% \begin{align*}
% \mathbb{E}\left[X^r\right] =&\ \frac{2C_{1,b}}{\pi i} \int_{c-i \cdot\infty}^{c+i \cdot\infty} \frac{\Gamma(s)}{\pi^s} \sum_{k=0}^\infty \frac{(-1)^k}{(2k+1)^{s}} \sum_{n=1}^\infty \frac{1}{n^{s-r}}\, \mathrm{d}s \\
% =&\ \frac{2C_{1,b}}{\pi i} \int_{c-i \cdot\infty}^{c+i \cdot\infty} \underbrace{\pi^{-s} \cdot \beta(s) \cdot \Gamma(s) \cdot \zeta(s - r)}_{=:\, \Xi(s)} \, \mathrm{d}s
% \end{align*}
% with Dirichlet's $\beta$, Euler's $\Gamma$ and Riemann's $\zeta$ functions.\\

% \textbf{Step 4:} We want to solve the integral by the method of residues by shifting the vertical line of integration to the left. Therefore we consider the analytic continuation of $\beta$ and the meromorphic continuations of $\Gamma$, and $\zeta$. In the representation $\Xi(s)$ of the integrand, the simple poles of $\Gamma$ and the trivial zeros of $\zeta(\cdot  - r)$ cancel each other out for even non-positive integers $s\in \{2-2n\}_{n\in \mathbb{N}} \subseteq \{r-2n\}_{n\in \mathbb{N}}$. In addition, the simple poles of $\Gamma$ and the simple zeros of $\beta$ cancel each other out for odd negative integers $s \in \{1-2n\}_{n \in \mathbb{N}}$. Hence, all $\Gamma$-poles in $s\in\{-n\}_{n\in \mathbb{N}_0}$ are canceled out and only the simple pole at $s = \varrho = r+1$ from $\zeta(\cdot - r)$ remains with residue
% $$
% \operatorname{Res}_{\Xi}(\varrho = r+1) = \pi^{-r-1} \cdot \beta(r+1) \cdot \Gamma(r+1) \cdot 1 = \frac{r!}{\pi^{r+1}} \beta(r+1) = \frac{(-1)^{\frac{r}{2}} \cdot E_r}{4 \cdot 2^r},
% $$
% since the identity
% $$
% \beta(r+1)=\frac{(-1)^{\frac{r}{2}} E_{r}}{2\cdot r!}\left(\frac\pi2\right)^{r+1}
% $$
% with Euler's numbers $E_r$ holds for even $r \in \mathbb{N}$.\\

% \textbf{Step 5:} We introduce a new contour for the integral of $\Xi$. Let $c=r+1+\varepsilon$ for some $\varepsilon>0$ and let $R= R_{1/2,c}(\pm T)$ be the rectangle with vertices $c \pm i\cdot T$ and $1/2 \pm i \cdot T$. Then,by the residue theorem, one has
% $$
% \frac{2 C_{1,b}}{\pi i}\int_{R} \Xi(s) \, \mathrm{d}s = 4 C_{1,b} \sum_{\varrho \in R} \operatorname{Res}_\Xi(\varrho) = C_{1,b} \frac{(-1)^{\frac{r}{2}}E_r}{2^r},
% $$
% while
% \begin{align*}
% &\ \frac{2 C_{1,b}}{\pi i}\int_{c-i\cdot T}^{c+i\cdot T} \Xi(s) \mathrm{d}s - C_{1,b} \frac{(-1)^{\frac{r}{2}}E_r}{2^r}\\
% =&\ \frac{2 C_{1,b}}{\pi i} \left[\underbrace{\int_{1/2 + i\cdot T}^{c + i\cdot T} \Xi(s) \mathrm{d}s}_{=: I_1} + \underbrace{\int_{1/2-i\cdot T}^{1/2+i\cdot T} \Xi(s) \mathrm{d}s}_{=: I_2} + \underbrace{\int_{c- i \cdot T}^{1/2 - i \cdot T} \Xi(s) \mathrm{d}s}_{=: I_3}\right].
% \end{align*}

% \textbf{Step 6:} We show that the integrals over the horizontal edges vanish, i.\,e.  $\lim_{T \rightarrow \infty} I_1 = \lim_{T \rightarrow \infty} I_3 = 0$. One has
% $$
% \lvert I_1\rvert, \lvert I_3\rvert \le \int_{1/2 \pm i\cdot T}^{c \pm i\cdot T} \lvert \pi^{-s} \beta(s) \Gamma(s) \zeta(s-r) \rvert \mathrm{d}s.
% $$
% Now, the Dirichlet's beta function can be written as a Dirichlet $L$-function, viz. $\beta(s) = L(s,\chi_4)$ for the alternating Dirichlet character of period 4. It exhibits at most polynomial growth. A classical result is the asymptotics $\beta(\sigma + i \cdot t) \in O\left((4\lvert t\rvert + 4)^{(1-\sigma)/2}\right) (\lvert t \rvert \rightarrow \infty)$ for $\sigma \in (0,1)$, cf. [Davenport, H. (1963). On Dirichlet's L-Functions. J. London Math. Soc. (3), 13: 524-536. Cited in: Francis, F. (2022). Explicit sub-convexity estimates for Dirichlet L-functions, https://arxiv.org/pdf/2206.11112]. Hence, for large t, $\beta$ grows at most at the rate of a square root, uniformly in $\sigma \in (0,1)$. For $\sigma > 1$, the series of the beta function converges absolutely, i.\,e. $\lvert \beta(s) \rvert = \lvert \sum_{k=0}^\infty (-1)^k/(2k+1)^s \rvert \le \sum_{k=0}^\infty (2k+1)^{-\sigma} \le \sum_{n=1}^\infty n^{-\sigma} = \zeta(\sigma) < \infty$. The singleton $\sigma = 1$ does not contribute to the integral. Furthermore, the functional equation of the zeta function, viz.
% $$\zeta(s) = 2^s \pi^{s-1} \sin\left(s \cdot \frac{\pi}{2}\right) \Gamma(1-s) \zeta(1-s)$$
% leads to
% $\lvert \zeta(s-r)\rvert = 2^{\sigma-r} \pi^{\sigma-r-1} \left\lvert \sin\left((s-r)\cdot \frac{\pi}{2}\right)\right\rvert \cdot \lvert \Gamma(1+r-s)\rvert \cdot \zeta(1+r-\sigma)$. Every term is well-behaved, i.\,e. uniformly bounded on $\sigma \in [1/2,1+r+\varepsilon]$, except for the zeta term as $\sigma$ tends to $r$. But in this case, the simple zero of the sine term cancels out the singularity of the zeta term. The singularity from the gamma term at $\sigma = 1+r$ was already canceled by the $\zeta$-factor in $\Xi$. Now, Stirling's formula results in 
% $$\lvert \Xi(\sigma + i \cdot T)\rvert \in O\left(T^{\sigma/2} e^{-T\pi/2}\right)(T \rightarrow \infty).$$
% Therefore, $\sup_{\sigma \in [1/2,1+r+\varepsilon]}\lvert \Xi(\sigma + i \cdot T) \rvert \rightarrow 0$ as $T \rightarrow \infty$ and $I_1$ and $I_3$ uniformly vanish in the limit.\\ 

% \textbf{Step 7:} Now, we show that $\lim_{T \rightarrow \infty} I_2$ exists. We recall classical results for the asymptotics of $\beta(s)$, represented as the Dirichlet L-function $L(s,\chi_4)$ and Riemann's $\zeta(s)$ on the important critical line $\sigma=1/2$. It is well known that $\beta$ and $\zeta$ exhibit at most polynomial growth as $\lvert t \rvert \rightarrow \infty$, cf. [Davenport, H. (1963). On Dirichlet's L-Functions. J. London Math. Soc. (3), 13: 524-536. Cited in: Francis, F. (2022). Explicit sub-convexity estimates for Dirichlet L-functions, https://arxiv.org/pdf/2206.11112] which yields $\beta(1/2 + i\cdot t) \in O(4(\lvert t \rvert)+1)^{\frac{1}{4}}$. Further, cf. [Bourgain, J. (2016). Decoupling, exponential sums and the Riemann zeta function. https://arxiv.org/pdf/1408.5794] shows $\zeta(1/2 + i \cdot t) \in O(\lvert t\rvert^{\frac{13}{84}+\varepsilon})$ for any $\delta > 0$. Together with Stirling's formula, those facts result in
% $$\Xi\left(\frac{1}{2} \pm i \cdot T\right) \in O\left(T^{\frac{17}{42} + \delta} e^{-T \pi/2}\right) (T\rightarrow\infty).$$
% Hence, due to the exponential decay, the expression 
% $$\int_{1/2-i\cdot \infty}^{1/2 + i \cdot \infty} \beta(s) \cdot \Gamma(s) \cdot \zeta(s-r)/\pi^s \mathrm{d}s = \int_{1/2-i\cdot \infty}^{1/2 + i \cdot \infty} \Xi(s) \mathrm{d}s$$
% exists, i.\,e. it is finite and unique.\\

% \textbf{Step 8:} From step 5 we deduce
% $$\int_{1/2-i\cdot \infty}^{1/2 + i \cdot \infty} \Xi(s) \mathrm{d}s = \frac{\pi i}{2C_{1,b}}\mathbb{E}\left[X^r\right] - \frac{(-1)^{\frac{r}{2}} i\pi \cdot E_r}{2^{r+1}}.$$
% Hence, we only have to evaluate $\mathbb{E}\left[X^r\right]$ for achieving the final result. Therefore, we use Poisson's summation formula. We have
% $$
% \mathbb{E}[X^r] = C_{1,b} \sum_{k \in \mathbb{Z}}^\infty k^r\cdot \operatorname{sech}\left(k \pi\right)
% $$
% Define the function $s(x) = C_{1,b} \cdot x^r \cdot \mathrm{sech}(x \pi)$ for every $r \in 2 \mathbb{N}$ with its Fourier transform
% \begin{align*}
% \mathcal{F}[s](\xi) =&\ \mathcal{F}[x \mapsto C_{1,b} \cdot x^r \cdot \mathrm{sech}(\pi x)] = C_{1,b} \int_{-\infty}^\infty s(x) e^{-i\cdot 2\pi x \xi} \mathrm{d}x \\
% =&\ C_{1,b} \cdot 2^{1-2r} \left(\sum_{k=1}^{r/2} (-1)^k T(r+1,k) \cosh (2k\pi \xi) + \frac{1}{2} T\left(r+1, \frac{r}{2}+1\right) \right) \mathrm{sech}^{r+1}(\pi \xi),
% \end{align*}
% with Euler's numbers $T(n,k)$ of type B, cf. https://oeis.org/A060187. Now, Poisson's summation formula states
% $$\sum_{k\in \mathbb{Z}}s(k) = \sum_{l \in \mathbb{Z}}\mathcal{F}[s](l).$$
% This results in
% $$\E[X^r] = C_{1,b} \cdot 2^{1-2r} \sum_{l = 1}^\infty \left(\sum_{k=1}^{r/2} (-1)^k T(r+1,k) \cosh (2k\pi l) + \frac{1}{2} T\left(r+1, \frac{r}{2}+1\right) \right) \mathrm{sech}^{r+1}(\pi l).$$

% \textbf{Step 9:} We evaluate the first identity in step 8. We get
% $$\int_{1/2-i\cdot \infty}^{1/2 + i \cdot \infty} \Xi(s) \mathrm{d}s = -\frac{\pi i}{16}\left(\sum_{k\in \mathbb{Z}} \cosh(2 \pi \xi)-3)\mathrm{sech}^3\left(\pi \xi \right)\right) - \frac{\pi i}{8}$$
% Finally,
% $$\cos\left(\frac{\pi}{4}\right) = \frac{1}{\sqrt{2}}.$$



% We show that the integrals over the horizontal edges, i.\,e. $I_1$ and $I_3$ vanish as $T \rightarrow \infty$. Stirling's formula (\ref{Stirlings_formula}) yields
% $$
% \left|\Gamma\left(\frac{1}{2} \pm i\cdot T\right)\right| \in O\left( e^{-|T| \cdot \pi/2}\right) \ (|T| \rightarrow \infty).
% $$
% Further
% $$
% \beta(\sigma \pm i \cdot T) \in O\left(|T|^A\right), \ \zeta(\sigma \pm i \cdot T) \in O\left(|T|^B\right)\ (|T| \rightarrow \infty)
% $$
% with at most polynomial growth rates $A,B$. We deduce
% $$
% |\Xi(\sigma \pm i\cdot T)| \in O\left(|T|^{A+B+\sigma-1/2} \cdot e^{-|T| \cdot \pi/2}\right) \ (|T| \rightarrow \infty)
% $$
% Therefore
% $$
% \sup_{\sigma \in [-N-1/2,c]} |\Xi(\sigma \pm i \cdot T)| \rightarrow 0 \ (|T|\rightarrow \infty).
% $$
% Since the horizontal edges have finite length, we deduce
% $$
% I_{\pm T} = \int_{c\pm i\cdot T}^{-N-1/2 \pm i\cdot T} \Xi(s) \, \mathrm{d}s \rightarrow 0 \ (|T| \rightarrow \infty).
% $$
% Next, the functional equation of Dirichlet's beta function, viz.
% $$
% \beta(s) = \left( \frac{\pi}{2} \right)^s \cdot \frac{\csc\left(s\cdot \frac{\pi}{2}\right)}{\Gamma(s)} \cdot \beta(1-s).
% $$
% yields
% \begin{align*}
% I_N &= \int_{-N-1/2-i\cdot T}^{-N-1/2+i\cdot T} \left(\frac{2}{\pi}\right)^{s} \beta(s) \cdot \Gamma(s) \cdot \zeta(s - r) \, \mathrm{d}s\\
% &= \int_{-N-1/2-i\cdot T}^{-N-1/2+i\cdot T} \csc\left(s\cdot \frac{\pi}{2}\right) \cdot \beta(1-s) \cdot \zeta(s-r) \, \mathrm{d}s.
% \end{align*}
% The cosecant term has simple poles at $s=-2k$ for $k\in \mathbb{N}_0$. Near that poles, choosing $\varepsilon << 1$, together with $s=-2k+\varepsilon \ \Leftrightarrow \varepsilon = s + 2k$ and the addition theorem for the sine result in
% \begin{align*}
% \sin\left(s\cdot \frac{\pi}{2}\right) &= \sin\left(-k\pi + \varepsilon \cdot \frac{\pi}{2}\right) = \sin(-k\pi)\cdot \cos\left(\varepsilon \cdot \frac{\pi}{2}\right) + \cos(-k\pi)\cdot \sin\left(\varepsilon \cdot \frac{\pi}{2}\right)\\
% &= (-1)^k \cdot \sin\left(\varepsilon \cdot \frac{\pi}{2}\right) \sim (-1)^{k} \cdot \frac{\pi}{2} \cdot \varepsilon \ (\varepsilon \downarrow 0).
% \end{align*}
% Therefore
% $$
% \operatorname{Res}_\Xi(-2k) = (-1)^k \cdot \frac{2}{\pi} \cdot \beta(1+2k)\cdot \zeta(-2k-r).
% $$
% And due to our choice $s>r+1$, we get an additional zeta pole with
% $$
% \operatorname{Res}_\Xi(s=r+1) = \frac{\beta(-r)}{\sin\left((r+1)\cdot \frac{\pi}{2}\right)}
% $$
% contributes to the integral. Hence, by shifting the left edge far to the left, we get
% $$
% \sum_{k=0}^\infty \frac{2}{\pi} (-1)^k \beta(1+2k) \cdot 
% $$




% we show $I_N \rightarrow 0$ uniformly in $t \in [-T,T]$ as $N \rightarrow \infty$. We apply Euler's reflection formula
% $$
% \Gamma(s) \cdot \Gamma(1-s) = \frac{\pi}{\sin(\pi \cdot s)}, \ s \not\in \mathbb{Z}
% $$
% to $s= -N -1/2 +i\cdot t \ \Rightarrow \ 1-s = N + 3/2 - i \cdot t$ results in
% $$
% \Gamma(-N-1/2 + i \cdot t) = \frac{\pi}{\sin(\pi(-N-1/2 + i \cdot t))) \cdot \Gamma(N + 3/2 - i \cdot t)}.
% $$
% Now the identity
% $$
% \sin(\sigma + i\cdot t) = \sin(\sigma) \cdot \cosh(t) + i \cdot \cos(\sigma) \cdot \sinh(t) 
% $$
% yields
% $$
% |\sin(\pi \cdot s)| = \left|(-1)^{N+1} \cosh(\pi \cdot t)\right| = \cosh(\pi \cdot t) \in O(e^{|t|\pi}/2) \ (t \rightarrow \infty).
% $$
% Hence,
% $$
% |\Gamma(s)| \lesssim \frac{e^{-|t|\pi}}{|\Gamma(N+3/2-i\cdot t)|}.
% $$
% Now, Stirling's formula in vertical strips yields
% $$
% |\Gamma(N+3/2 - i \cdot t)| \asymp N^{N+1}\cdot e^{-N-3/2},
% $$
% uniformly in $t$, since
% $$
% \Gamma(s) \in \sqrt{2\pi} s^{s-1/2} e^{-s}\left(1+O\left(\frac{1}{|s|}\right)\right)
% $$
% for $|s| \rightarrow \infty$ in sectors $|\arg(s)|\leq \pi - \delta$ for some small $\delta > 0$ which is fulfilled for $|t|>0$. Hence,
% $$
% \Gamma(-N-1/2 + i \cdot t) \in O\left(N^{-N-1}\cdot e^{N-3/2}\right).
% $$
% Further, the reflection formula for Riemann's zeta function yields
% $$
% \zeta(r-s) = 2^{s-r} \cdot \pi^{s-r-1} \cdot \sin \left(\frac{\pi(s-r)}{2}\right) \cdot \Gamma(1-(s-r)) \cdot \zeta(1-(s-r)).
% $$
% Now, $1-(s-r) = N + 3/2 +r - i\cdot t$ yields $\zeta(N + 3/2 + r - i \cdot t) \rightarrow 1 \ (N \rightarrow \infty)$. A further application of Stirling's formula shows
% $$
% |\Gamma(N + 3/2 +r - i\cdot t)| \in O\left(N^{N+r+1} \cdot e^{-N}\right) \ (N \rightarrow \infty).
% $$
% Hence,
% $$
% |\zeta(-N-1/2 + i\cdot t)| \in O\left((2\pi)^{-N-1/2-r} \cdot N^{N+r+1} \cdot e^{-N}\right)\ (N \rightarrow \infty).
% $$
% Moreover, the functional equation
% $$
% \beta(s) = \left( \frac{\pi}{2}\right)^{s} \frac{1}{\sin(s\cdot \pi/2)} \frac{1}{\Gamma(s)}\cdot \beta(1-s)
% $$
% with $\beta(1+N+1/2 - i\cdot t) \rightarrow 1 \ (N \rightarrow \infty)$ and again double application of Stirling's formula on vertical stripes yield
% $$
% \left|\frac{\Gamma((1-s)/2)}{\Gamma(1+s)/2}\right| \in O\left(N^{N+1/2}\cdot e^{-N}\right)\ (N \rightarrow \infty).
% $$
% Therefore,
% $$
% |\beta(-N-1/2 + i \cdot t)| \in O\left(\left(\frac{4}{\pi}\right)^{N+1} \cdot N^{N+1/2}\cdot e^{-N}\right)\ (N \rightarrow \infty).
% $$

% Thus, the residue theorem yields the following explicit formula
% $$
% \mu_r = \frac{4}{(1+\sqrt{2})G - 1} \cdot \left( r! \cdot \beta(1+r) \cdot \left(\frac{2}{\pi}\right)^{1+r} + \sum_{n=0}^\infty \frac{(-1)^n}{n!} \cdot \beta(-n) \cdot \zeta(-n-r) \cdot \left(\frac{2}{\pi}\right)^{-n} \right)
% $$
% for the moments of order $r$ of our discrete hyperbolic secant distribution. We simplify the formula by observing its behavior for even parities of $r=2j$ and $m=2k$ as well as odd parities of $r=2j+1$ and $n=2k+1$ with $j \in \mathbb{N}$ and $k \in \mathbb{N}_0$. For Dirichlet's  beta function, we observe
% $$
% \beta(2j+1) = \frac{(-1)^j E_{2j}}{2(2j)!}\cdot \left( \frac{\pi}{2} \right)^{2j+1}, \quad \beta(-2j) = \begin{cases}
%     \frac{1}{2}E_{2j}, & j\geq1,\\
%     \frac{1}{2}, & j=0
% \end{cases}
% $$


% and by splitting the series in even indices $n=$ and odd indices. We denote the $n$-th \textit{Euler number} by $E_n$. If $r=2j$ is even, then  we get
% $$
% \beta(2j+1) = \frac{(-1)^j E_{2j}}{2(2j)!}\cdot \left( \frac{\pi}{2} \right)^{2j+1}, \quad \beta(-2j) = \begin{cases}
%     \frac{1}{2}E_{2j}, & j\geq1,\\
%     \frac{1}{2}, & j=0
% \end{cases}
% $$
% and the trivial zeros of Riemann's zeta function yield
% $$
% \zeta(-m) = 0 \quad \text{ for } m \in 2 \mathbb{N}.
% $$
% If $r=2j+1$ is odd, then there is no closed form for $\beta(2j+2)$ and $\beta(2)$ is called \textit{Catalan's constant}. Further,
% $$
% \beta(-(2j+1)) = 0, \quad \zeta(-n-(2j+1)) = \begin{cases}
%     0, & n \text{ odd},\\
%     \zeta(-n-2j-1), & \text{else}.
% \end{cases}
% $$
% We split the series into odd and even parts and deduce
% $$
% \mu_{r=2j} =  \frac{2}{(1+\sqrt{2})G - 1}\left( \right)
% $$
% \end{proof}

% \begin{theorem}
%     Let $X\sim\mathsf{Sech}_{3,u}$ be a discrete random variable which is distributed according to the \textit{unilateral hyperbolic secant distribution of the third kind}. Let $r\in\N$ be a fixed positive integer. Then the moments of order $r$ of $X$ can be represented by
%     \begin{align*}
%         \E\left[X^r\right] &= \frac{4}{(1+\sqrt{2})G + 1} \left[r! \beta(r+1)\left(\frac{2}{\pi}\right)^{r+1} + \sum_{n=0}^\infty \frac{(-1)^n}{n!} \left(\frac{2}{\pi}\right)^{-n} \zeta(-r-n) \beta(-n)\right].
%     \end{align*}
%     For even $r$, this results in
%     $$
%         \E\left[X^{r}\right] = \frac{2 \cdot (-1)^{r/2} E_{r}}{(1+\sqrt{2})G-1}
%     $$
%     with the \textit{Euler numbers} $E_r$.
% \end{theorem}

% \begin{proof}
% Let $r$ be a fixed positive integer. We calculate the $r$-th moments $\mu_r$ of a random variable $X \sim \mathsf{Sech}_{3,u}$. By definition, one has
% $$
% \mu_r := \mathbb{E}\left[X^r\right] = \frac{2}{(1+\sqrt{2})G + 1} \sum_{n=0}^\infty n^r\cdot \operatorname{sech}\left(n \cdot \frac{\pi}{2}\right).
% $$
% Let $x > 0$. Then one has $|\exp(-2x)| < 1$. Hence, the geometric series yields the identities
% $$
% \operatorname{sech}(x) = \frac{2e^{-x}}{1+e^{-2x}} = 2e^{-x} \sum_{k=0}^\infty (-1)^k e^{-2kx} = 2 \sum_{k=0}^\infty (-1)^k e^{-(2k+1)x}.
% $$
% Now, the gamma function motivates the introduction of the \textit{Mellin integral transform} $\mathcal{M}$. For every function $f \in L^1_{\mathrm{loc}}((0,\infty), \mathbb{C})$ and parameter $s \in \mathbb{C}$, it is defined as
% $$
% \mathcal{M}\left\{f\right\}(s) = \varphi(s)=\int_0^\infty x^{s-1} f(x) \, \mathrm{d}x.
% $$
% The integral converges absolutely in a strip $\mathbb{S}(a,b) := \{\sigma \in \mathbb{C}: a < \Re(\sigma) < b\}$ for $a < b$. Now, $[x\mapsto e^{-x}] \in L^1_{\mathrm{loc}}((0,\infty), \mathbb{C})$ and we observe
% $$
% \mathcal{M}\{\left[x \mapsto e^{-x}\right]\} = \int_0^\infty x^{s-1}e^{-x} \, \mathrm{d}x = \Gamma(s).
% $$
% The gamma function is holomorphic on the open right half-plane. Therefore, it is in particular analytic on $\mathbb{S}(0,b)$ for every finite $b>0$. Further, Stirling's formula
% \begin{equation}\label{Stirlings_formula}
% |\Gamma(c + i \cdot t)| \sim \sqrt{2\pi} |t|^{c-1/2} e^{-|t|\cdot \pi/2}(1+O(1/|t|)) \ (|t| \rightarrow \infty)
% \end{equation}
% holds uniformly for any $c=\Re(s) \in \mathbb{S}(0,b)$. Thus,
% $$
% \int_{c-i \cdot\infty}^{c+i \cdot\infty} |\Gamma(s)| \, \mathrm{d}s < \infty.
% $$
% Hence, for every $c>0$, the Mellin inversion theorem yields the \textit{Cahen-Mellin integral}
% $$
% e^{-x} = \mathcal{M}^{-1}\left\{\Gamma\right\}(x) = \frac{1}{2\pi i} \int_{c-i \cdot\infty}^{c+i \cdot\infty} \Gamma(s) \cdot x^{-s} \, \mathrm{d}s.
% $$
% This yields
% $$
% \operatorname{sech}(x) = 2 \sum_{k=0}^\infty (-1)^k e^{-(2k+1)x} = \sum_{k=0}^\infty (-1)^k \frac{2}{2\pi i} \int_{c-i \cdot\infty}^{c+i \cdot\infty} \Gamma(s) \cdot (2k+1)^{-s} \cdot x^{-s}\, \mathrm{d}s.
% $$

% Therefore
% \begin{align*}
% \mu_r &= \frac{2}{(1+\sqrt{2})G + 1} \sum_{n=0}^\infty n^r\cdot \operatorname{sech}\left(n \cdot \frac{\pi}{2}\right)\\
% &= \frac{2}{(1+\sqrt{2})G + 1} \sum_{n=1}^\infty n^r \sum_{k=0}^\infty (-1)^k \frac{2}{2\pi i} \int_{c-i \cdot\infty}^{c+i \cdot\infty} \frac{\Gamma(s)}{(2k+1)^{s}} \cdot \left(n \cdot \frac{\pi}{2}\right)^{-s}\, \mathrm{d}t.
% \end{align*}
% Now, we want to interchange the limits in the sums and the integral. Choose $c>r+1$. On the line $\Re(s)=c$ one has
% $$
% \left|(-1)^k \cdot \Gamma(s) \cdot (2k+1)^{-s} \cdot n^{r-s} \cdot \left(\frac{\pi}{2}\right)^{-s}\right| = \left(\frac{2}{\pi}\right)^{c} \cdot |\Gamma(c)| \cdot (2k+1)^{-c} \cdot n^{r-c}
% $$
% Stirling's formula (\ref{Stirlings_formula}) and $c>r+1$ yield
% $$
% \int_{c-i \cdot\infty}^{c+i \cdot\infty} |\Gamma(s)|\, \mathrm{d}s < \infty, \quad \sum_{k=0}^\infty (2k+1)^{-c} < \infty, \quad \sum_{n=1}^\infty n^{r-c} < \infty.
% $$
% Thus, Tonelli's theorem and the functional equation of Dirichlet's beta function, viz.
% $$
% \beta(s) = \left( \frac{\pi}{2} \right)^s \cdot \frac{\csc\left(s\cdot \frac{\pi}{2}\right)}{\Gamma(s)} \cdot \beta(1-s).
% $$
% yield
% \begin{align*}
% \mu_r=&\ \frac{4}{(1+\sqrt{2})G + 1}  \cdot \frac{1}{2\pi i} \int_{c-i \cdot\infty}^{c+i \cdot\infty} \left(\frac{2}{\pi}\right)^{s} \Gamma(s) \sum_{k=0}^\infty \frac{(-1)^k}{(2k+1)^{s}} \sum_{n=1}^\infty \frac{1}{n^{s - r}}\, \mathrm{d}s \\
% =&\ \frac{4}{(1+\sqrt{2})G + 1} \cdot \frac{1}{2\pi i} \int_{c-i \cdot\infty}^{c+i \cdot\infty} \left(\frac{2}{\pi}\right)^{s} \beta(s) \cdot \Gamma(s) \cdot \zeta(s - r) \, \mathrm{d}s \\
% =&\ \frac{4}{(1+\sqrt{2})G + 1} \cdot \frac{1}{2\pi i} \int_{c-i \cdot\infty}^{c+i \cdot\infty} \underbrace{\csc\left(s \cdot \frac{\pi}{2}\right) \cdot \beta(1-s) \cdot \zeta(s-r)}_{=: \Xi(s)} \, \mathrm{d}s 
% \end{align*}
% with \textit{Dirichlet's beta function} $\beta$ and \textit{Riemann's zeta function} $\zeta$. We want to solve the integral by the method of residues. Therefore we consider the meromorphic continuations of $\beta$, and $\zeta$. Our function $\Xi$ has poles $\varrho \in \{2z\}_{z\in \mathbb{Z}} \, \cup \, \{1+r\}$ with residues
% $$
% \operatorname{Res}_\Xi(\varrho=z) = \left(\frac{2}{\pi}\right)^{z} \beta(-n) \cdot \frac{(-1)^z}{!} \cdot \zeta(z-r) 
% $$
% and
% $$
% \operatorname{Res}_\Xi(\varrho=1+r) = \left(\frac{2}{\pi}\right)^{1+r} \beta(1+r) \cdot r!. 
% $$
% We introduce a new contour for the path integral. Let $R_{N,T}$ be the rectangle with vertices $c \pm i\cdot T$ and $-N - 1/2 \pm i \cdot T$. Then,by the residue theorem, one has
% $$
% \int_{R_{N,T}} \Xi(s) \, \mathrm{d}s = 2\pi i \sum_{\varrho \in R_{N,T}} \operatorname{Res}_\Xi(\varrho),
% $$
% while
% $$
% \int_{R_{N,T}} \Xi(s) \, \mathrm{d}s = \int_{c-i\cdot T}^{c+i\cdot T} \Xi(s) \, \mathrm{d}s + \underbrace{\int_{c\pm i\cdot T}^{-N-1/2 \pm i\cdot T} \Xi(s) \, \mathrm{d}s}_{=: I_{\pm T}} + \underbrace{\int_{-N-1/2-i\cdot T}^{-N-1/2+i\cdot T} \Xi(s) \, \mathrm{d}s}_{=: I_N}.
% $$
% We show that the horizontal edges $I_{\pm T} \rightarrow 0$ uniformly in $\sigma \in [-N-1/2,c]$ as $|T| \rightarrow \infty$. Stirling's formula (\ref{Stirlings_formula}) yields
% $$
% |\Gamma(\sigma \pm i\cdot T)| \in O\left(|T|^{\sigma-1/2} \cdot e^{-|T| \cdot \pi/2}\right) \ (|T| \rightarrow \infty).
% $$
% Further
% $$
% \beta(\sigma \pm i \cdot T) \in O\left(|T|^A\right), \ \zeta(\sigma \pm i \cdot T) \in O\left(|T|^B\right)\ (|T| \rightarrow \infty)
% $$
% and for some polynomial growth rates $A,B$. We deduce
% $$
% |\Xi(\sigma \pm i\cdot T)| \in O\left(|T|^{A+B+\sigma-1/2} \cdot e^{-|T| \cdot \pi/2}\right) \ (|T| \rightarrow \infty)
% $$
% Therefore
% $$
% \sup_{\sigma \in [-N-1/2,c]} |\Xi(\sigma \pm i \cdot T)| \rightarrow 0 \ (|T|\rightarrow \infty).
% $$
% Since the horizontal edges have finite length, we deduce
% $$
% I_{\pm T} = \int_{c\pm i\cdot T}^{-N-1/2 \pm i\cdot T} \Xi(s) \, \mathrm{d}s \rightarrow 0 \ (|T| \rightarrow \infty).
% $$
% Next, the functional equation of Dirichlet's beta function, viz.
% $$
% \beta(s) = \left( \frac{\pi}{2} \right)^s \cdot \frac{\csc\left(s\cdot \frac{\pi}{2}\right)}{\Gamma(s)} \cdot \beta(1-s).
% $$
% yields
% \begin{align*}
% I_N &= \int_{-N-1/2-i\cdot T}^{-N-1/2+i\cdot T} \left(\frac{2}{\pi}\right)^{s} \beta(s) \cdot \Gamma(s) \cdot \zeta(s - r) \, \mathrm{d}s\\
% &= \int_{-N-1/2-i\cdot T}^{-N-1/2+i\cdot T} \csc\left(s\cdot \frac{\pi}{2}\right) \cdot \beta(1-s) \cdot \zeta(s-r) \, \mathrm{d}s.
% \end{align*}
% The cosecant term has simple poles at $s=-2k$ for $k\in \mathbb{N}_0$. Near that poles, choosing $\varepsilon << 1$, together with $s=-2k+\varepsilon \ \Leftrightarrow \varepsilon = s + 2k$ and the addition theorem for the sine result in
% \begin{align*}
% \sin\left(s\cdot \frac{\pi}{2}\right) &= \sin\left(-k\pi + \varepsilon \cdot \frac{\pi}{2}\right) = \sin(-k\pi)\cdot \cos\left(\varepsilon \cdot \frac{\pi}{2}\right) + \cos(-k\pi)\cdot \sin\left(\varepsilon \cdot \frac{\pi}{2}\right)\\
% &= (-1)^k \cdot \sin\left(\varepsilon \cdot \frac{\pi}{2}\right) \sim (-1)^{k} \cdot \frac{\pi}{2} \cdot \varepsilon \ (\varepsilon \downarrow 0).
% \end{align*}
% Therefore
% $$
% \operatorname{Res}_\Xi(-2k) = (-1)^k \cdot \frac{2}{\pi} \cdot \beta(1+2k)\cdot \zeta(-2k-r).
% $$
% And due to our choice $s>r+1$, we get an additional zeta pole with
% $$
% \operatorname{Res}_\Xi(s=r+1) = \frac{\beta(-r)}{\sin\left((r+1)\cdot \frac{\pi}{2}\right)}
% $$
% contributes to the integral. Hence, by shifting the left edge far to the left, we get
% $$
% \sum_{k=0}^\infty \frac{2}{\pi} (-1)^k \beta(1+2k) \cdot 
% $$




% we show $I_N \rightarrow 0$ uniformly in $t \in [-T,T]$ as $N \rightarrow \infty$. We apply Euler's reflection formula
% $$
% \Gamma(s) \cdot \Gamma(1-s) = \frac{\pi}{\sin(\pi \cdot s)}, \ s \not\in \mathbb{Z}
% $$
% to $s= -N -1/2 +i\cdot t \ \Rightarrow \ 1-s = N + 3/2 - i \cdot t$ results in
% $$
% \Gamma(-N-1/2 + i \cdot t) = \frac{\pi}{\sin(\pi(-N-1/2 + i \cdot t))) \cdot \Gamma(N + 3/2 - i \cdot t)}.
% $$
% Now the identity
% $$
% \sin(\sigma + i\cdot t) = \sin(\sigma) \cdot \cosh(t) + i \cdot \cos(\sigma) \cdot \sinh(t) 
% $$
% yields
% $$
% |\sin(\pi \cdot s)| = \left|(-1)^{N+1} \cosh(\pi \cdot t)\right| = \cosh(\pi \cdot t) \in O(e^{|t|\pi}/2) \ (t \rightarrow \infty).
% $$
% Hence,
% $$
% |\Gamma(s)| \lesssim \frac{e^{-|t|\pi}}{|\Gamma(N+3/2-i\cdot t)|}.
% $$
% Now, Stirling's formula in vertical strips yields
% $$
% |\Gamma(N+3/2 - i \cdot t)| \asymp N^{N+1}\cdot e^{-N-3/2},
% $$
% uniformly in $t$, since
% $$
% \Gamma(s) \in \sqrt{2\pi} s^{s-1/2} e^{-s}\left(1+O\left(\frac{1}{|s|}\right)\right)
% $$
% for $|s| \rightarrow \infty$ in sectors $|\arg(s)|\leq \pi - \delta$ for some small $\delta > 0$ which is fulfilled for $|t|>0$. Hence,
% $$
% \Gamma(-N-1/2 + i \cdot t) \in O\left(N^{-N-1}\cdot e^{N-3/2}\right).
% $$
% Further, the reflection formula for Riemann's zeta function yields
% $$
% \zeta(r-s) = 2^{s-r} \cdot \pi^{s-r-1} \cdot \sin \left(\frac{\pi(s-r)}{2}\right) \cdot \Gamma(1-(s-r)) \cdot \zeta(1-(s-r)).
% $$
% Now, $1-(s-r) = N + 3/2 +r - i\cdot t$ yields $\zeta(N + 3/2 + r - i \cdot t) \rightarrow 1 \ (N \rightarrow \infty)$. A further application of Stirling's formula shows
% $$
% |\Gamma(N + 3/2 +r - i\cdot t)| \in O\left(N^{N+r+1} \cdot e^{-N}\right) \ (N \rightarrow \infty).
% $$
% Hence,
% $$
% |\zeta(-N-1/2 + i\cdot t)| \in O\left((2\pi)^{-N-1/2-r} \cdot N^{N+r+1} \cdot e^{-N}\right)\ (N \rightarrow \infty).
% $$
% Moreover, the functional equation
% $$
% \beta(s) = \left( \frac{\pi}{2}\right)^{s} \frac{1}{\sin(s\cdot \pi/2)} \frac{1}{\Gamma(s)}\cdot \beta(1-s)
% $$
% with $\beta(1+N+1/2 - i\cdot t) \rightarrow 1 \ (N \rightarrow \infty)$ and again double application of Stirling's formula on vertical stripes yield
% $$
% \left|\frac{\Gamma((1-s)/2)}{\Gamma(1+s)/2}\right| \in O\left(N^{N+1/2}\cdot e^{-N}\right)\ (N \rightarrow \infty).
% $$
% Therefore,
% $$
% |\beta(-N-1/2 + i \cdot t)| \in O\left(\left(\frac{4}{\pi}\right)^{N+1} \cdot N^{N+1/2}\cdot e^{-N}\right)\ (N \rightarrow \infty).
% $$

% Thus, the residue theorem yields the following explicit formula
% $$
% \mu_r = \frac{4}{(1+\sqrt{2})G - 1} \cdot \left( r! \cdot \beta(1+r) \cdot \left(\frac{2}{\pi}\right)^{1+r} + \sum_{n=0}^\infty \frac{(-1)^n}{n!} \cdot \beta(-n) \cdot \zeta(-n-r) \cdot \left(\frac{2}{\pi}\right)^{-n} \right)
% $$
% for the moments of order $r$ of our discrete hyperbolic secant distribution. We simplify the formula by observing its behavior for even parities of $r=2j$ and $m=2k$ as well as odd parities of $r=2j+1$ and $n=2k+1$ with $j \in \mathbb{N}$ and $k \in \mathbb{N}_0$. For Dirichlet's  beta function, we observe
% $$
% \beta(2j+1) = \frac{(-1)^j E_{2j}}{2(2j)!}\cdot \left( \frac{\pi}{2} \right)^{2j+1}, \quad \beta(-2j) = \begin{cases}
%     \frac{1}{2}E_{2j}, & j\geq1,\\
%     \frac{1}{2}, & j=0
% \end{cases}
% $$


% and by splitting the series in even indices $n=$ and odd indices. We denote the $n$-th \textit{Euler number} by $E_n$. If $r=2j$ is even, then  we get
% $$
% \beta(2j+1) = \frac{(-1)^j E_{2j}}{2(2j)!}\cdot \left( \frac{\pi}{2} \right)^{2j+1}, \quad \beta(-2j) = \begin{cases}
%     \frac{1}{2}E_{2j}, & j\geq1,\\
%     \frac{1}{2}, & j=0
% \end{cases}
% $$
% and the trivial zeros of Riemann's zeta function yield
% $$
% \zeta(-m) = 0 \quad \text{ for } m \in 2 \mathbb{N}.
% $$
% If $r=2j+1$ is odd, then there is no closed form for $\beta(2j+2)$ and $\beta(2)$ is called \textit{Catalan's constant}. Further,
% $$
% \beta(-(2j+1)) = 0, \quad \zeta(-n-(2j+1)) = \begin{cases}
%     0, & n \text{ odd},\\
%     \zeta(-n-2j-1), & \text{else}.
% \end{cases}
% $$
% We split the series into odd and even parts and deduce
% $$
% \mu_{r=2j} =  \frac{2}{(1+\sqrt{2})G - 1}\left( \right)
% $$
% \end{proof}

% Letb $r\in \mathbb{N}$ be a positive, even number. For a complex parameter $s$, define Dirichlet's beta fucntion as $\beta(s) = \sum_{n=0}^\infty \frac{(-1)^n} {(2n+1)^s}$. Further, define the gamma function $\Gamma(s) = \int_0^\infty t^{s-1} e^{-t} \mathrm{d}t$ and the zeta function $\zeta(s)=\sum_{n=1}^\infty \frac{1}{n^s}$. Define $\Xi(s):=\beta(s)\cdot \Gamma(s) \cdot \zeta(s-r)$. Calculate $\int_{1/2-i\cdot \infty}^{1/2 + i \cdot \infty} \Xi(s) \mathrm{d}s$.

% Return an exact mathematical expression, you can involve numbers and special fucntions.



\end{document}
